\documentclass[11pt,a4paper]{article}

\usepackage[utf8]{inputenc}
\usepackage[T1]{fontenc}
\usepackage{babel}
\babelprovide[import,main]{ngerman}
\babelprovide[import]{english}
\usepackage{lmodern}
\usepackage{microtype}
\usepackage{geometry}
\usepackage{amsmath,amssymb,amsthm,mathtools}
\usepackage{booktabs,longtable,tabularx,array}
\usepackage{enumitem}
\usepackage{xcolor}
\usepackage{fancyhdr}
\usepackage{seqsplit}
\usepackage[numbers,sort&compress]{natbib}
\usepackage{hyperref}
\usepackage[nameinlink,noabbrev]{cleveref}

\geometry{left=25mm,right=25mm,top=25mm,bottom=29mm,headheight=14pt}
\setlength{\parindent}{0pt}
\setlength{\parskip}{0.58em}
\setlength{\emergencystretch}{3em}
\setlist[itemize]{leftmargin=1.45em,itemsep=0.2em,topsep=0.3em}
\setlist[enumerate]{leftmargin=1.7em,itemsep=0.25em,topsep=0.3em}

\definecolor{QBlue}{HTML}{17365D}
\definecolor{QLight}{HTML}{EDF3F8}
\definecolor{QGreen}{HTML}{EAF4EC}
\definecolor{QAmber}{HTML}{FFF4D6}
\definecolor{QRed}{HTML}{FBE9E7}
\definecolor{QGray}{HTML}{4A4A4A}

\hypersetup{
  colorlinks=true,
  linkcolor=QBlue,
  citecolor=QBlue,
  urlcolor=QBlue,
  pdftitle={Mandelbrot-Menge, Anschlussordnung, Physik und Retrokausalitaet},
  pdfauthor={Ingolf Lohmann},
  pdfsubject={Komplementbeweis, rekursive Wirkungsklassifikation, Dimensionsbruecke, Retrokausalitaet und Ontologie des Unterschieds},
  pdfkeywords={Mandelbrot-Menge, Komplement, Anschlussordnung, Fixpunkt, Quantisierbarkeit, Dimensionsanalyse, Raumzeit, Quantengravitation, Retrokausalitaet, Retrodiktion, No-Signalling, Ontologie des Unterschieds, QIK-VRT}
}

\pagestyle{fancy}
\fancyhf{}
\fancyhead[L]{\small\textcolor{QGray}{Mandelbrot, Anschlussordnung, Physik und Zeit}}
\fancyhead[R]{\small\textcolor{QGray}{Ingolf Lohmann}}
\fancyfoot[C]{\small\thepage}
\renewcommand{\headrulewidth}{0.3pt}
\renewcommand{\footrulewidth}{0pt}

\newtheorem{definition}{Definition}[section]
\newtheorem{satz}{Satz}[section]
\newtheorem{lemma}{Lemma}[section]
\newtheorem{korollar}{Korollar}[section]
\newtheorem{proposition}{Proposition}[section]
\theoremstyle{remark}
\newtheorem{bemerkung}{Bemerkung}[section]
\newtheorem{gegenmodell}{Gegenmodell}[section]
\crefname{definition}{Definition}{Definitionen}
\crefname{satz}{Satz}{Sätze}
\crefname{lemma}{Lemma}{Lemmata}
\crefname{korollar}{Korollar}{Korollare}
\crefname{proposition}{Proposition}{Propositionen}
\crefname{bemerkung}{Bemerkung}{Bemerkungen}
\crefname{gegenmodell}{Gegenmodell}{Gegenmodelle}

\newcommand{\M}{\mathcal{M}}
\newcommand{\K}{\mathcal{K}}
\newcommand{\C}{\mathbb{C}}
\newcommand{\R}{\mathbb{R}}
\newcommand{\N}{\mathbb{N}}
\newcommand{\eps}{\varepsilon}
\newcommand{\dd}{\mathop{}\!\mathrm{d}}
\newcommand{\qikvrt}{QIK--VRT}
\newcommand{\Fix}{\operatorname{Fix}}
\newcommand{\PASS}{\texttt{PASS}}
\newcommand{\CONTINUE}{\texttt{CONTINUE}}
\newcommand{\BLOCK}{\texttt{BLOCK}}
\newcommand{\hash}[1]{{\ttfamily\small\seqsplit{#1}}}
\newcommand{\doilink}[1]{\href{https://doi.org/#1}{\nolinkurl{#1}}}
\newcommand{\statusbox}[2]{%
  \par\medskip\noindent
  \colorbox{#1}{\parbox{\dimexpr\linewidth-2\fboxsep\relax}{#2}}%
  \par\medskip
}
\newcommand{\dotcup}{\mathbin{\mathaccent\cdot\cup}}

\begin{document}

\hypersetup{pageanchor=false}
\begin{titlepage}
  \thispagestyle{empty}
  \vspace*{16mm}
  {\color{QBlue}\rule{\textwidth}{1.4pt}}\\[1.2em]
  {\Huge\bfseries Mandelbrot-Menge,\\[0.2em]
  Anschlussordnung und\\[0.2em]
  physikalisches Modelluniversum\par}
  \vspace{0.8em}
  {\Large Komplementbeweis, rekursive Wirkungsklassifikation,\\
  dimensionsrichtige Raumzeitbrücke, Retrokausalität und der\\
  entscheidende Unterschied zur empirischen Quantengravitation\par}
  \vspace{1.2em}
  {\color{QBlue}\rule{\textwidth}{0.6pt}}

  \vfill
  {\large\textbf{Ingolf Lohmann}\par}
  \vspace{0.4em}
  {\normalsize Wissenschaftlich standardisierte und prämissentransparente Fassung\par}
  \vspace{0.3em}
  {\normalsize Version 3.0 -- 21. Juli 2026\par}

  \vfill
  \statusbox{QLight}{
    \textbf{Wissenschaftlicher Status.}
    Dieses Dokument enthält mathematische Sätze mit Beweisen, ausdrücklich
    definierte Modellbegriffe, physikalische Korrespondenzhypothesen,
    ontologische Interpretationen und eine normative Schlussfolgerung.
    Diese Aussagearten werden nicht miteinander gleichgesetzt.
  }

  \vfill
  {\small
    Dokumentationslizenz: CC BY-NC-ND 4.0, soweit nicht bei zitierten
    Fremdwerken anders angegeben. Mathematische Standardbegriffe und
    Gleichungen bleiben davon unberührt.
  }
\end{titlepage}
\hypersetup{pageanchor=true}
\setcounter{page}{1}

\selectlanguage{ngerman}

\begin{abstract}
Teil I beweist, dass die Mandelbrot-Menge \(\M\) und ihr relatives Komplement
exakt den zuvor festgelegten Grundbereich \(U\) ergeben:
\[
  U=(\M\cap U)\dotcup\bigl(U\setminus(\M\cap U)\bigr).
\]
Dieser Satz ist in klassischer Mengenlehre vollständig bewiesen. Wird \(U\)
als \emph{Modelluniversum} bezeichnet, so ergeben die eingeschränkte
Mandelbrot-Klasse \(\M_U\) und ihr relatives Komplement \(\K_U\) genau dieses
Modelluniversum. Eine Aussage über den physikalischen
Kosmos entsteht dagegen erst durch eine explizite Korrespondenzabbildung von
Raumzeitereignissen in den Modellbereich. Das Dokument definiert deshalb eine
raumzeitlich induzierte Komplementärmenge als Urbildpartition und trennt deren
mathematische Vollständigkeit von der offenen Frage ihrer physikalischen
Adäquatheit.

Ferner wird bewiesen, dass jeder beschränkte Ausschnitt eines
endlichdimensionalen euklidischen Modellraums bei jeder vorgegebenen positiven
Genauigkeit endlich diskret darstellbar ist. Diese
\emph{darstellungsbezogene Quantisierbarkeit} impliziert keine ontische
Diskretheit und keine fundamentale Quantisierung der Raumzeit; der
kontinuierliche Einheitswürfel \([0,1]^4\) ist ein unmittelbares Gegenmodell zu diesem
Fehlschluss. Gerade
\[
  \text{Quantisierbarkeit der Darstellung}
  \neq
  \text{Quantisierung des Dargestellten}
\]
ist der entscheidende Unterschied.

Die klassische Allgemeine Relativitätstheorie kann das Innere idealisierter
Schwarzer Löcher mathematisch darstellen, obwohl es für äußere Beobachter
kausal nicht einsehbar ist. Das ``geistige Auge'' bezeichnet hier präzise
theoretische Darstellbarkeit, nicht Messzugang. Die Planck-Skala ist eine aus
\(\hbar\), \(G\) und \(c\) gebildete charakteristische Skala und motiviert
Quantengravitation; sie ist weder als kleinste Länge noch als ontische
Raumzeitkörnung experimentell bewiesen. Abschließend wird gezeigt, in welchem
strengen Sinn diese Unterschiede an die Ontologie des universalen
Unterschieds und an ein universal formuliertes Schöpfungsprinzip anschließen,
ohne eine metaphysische Interpretation als bereits bestätigtes Naturgesetz
auszugeben.

Teil II überführt die statische Zerlegung in ein typisiertes Prozessmodell.
Die erweiterte Dynamik \(F(c,z)=(c,z^2+c)\) erzeugt Trajektorien; endliche
Gates unterscheiden \(\PASS\), \(\CONTINUE\) und \(\BLOCK\), ohne
aus bloßer Nichtflucht eine Mitgliedschaft zu folgern. Bewiesen werden die
Shift-Invarianz der exakten Klassifikation, die monotone Verfeinerung
zertifizierter Zustände und ein Faktorisierungskriterium für gate-erhaltende
Prozessabbildungen. Eine spätere terminale Klassifikation bestimmt die
Bedeutung einer früheren Spur neu, verändert aber keinen früheren Zustand und
begründet daher keine physikalische Retrokausalität. Der tragfähige Fixpunkt
liegt auf der Ebene der Klassifikation, nicht allgemein im Einzelorbit.

Teil III konstruiert die fehlende Dimensionsbrücke. Er trennt mathematische
Metrik, Lorentz-Metrik, Messzahl und Einheit; führt dimensionslose Variablen
durch explizite Referenzskalen ein; prüft die Dimensionen von
Einstein-Gleichung, Wirkung, Planck-Größen und Horizontentropie; und formuliert
notwendige Bedingungen für eine physikalische Korrespondenz. Damit wird
präzisiert, was logisch, mathematisch und dimensionsanalytisch folgt und
welche dynamischen Gleichungen, Observablen und Messdaten für eine Theorie der
Quantengravitation weiterhin fehlen.

Teil IV untersucht Retrokausalität als eigenständige physikalische und
erkenntnis\-theoretische Hypothesenfamilie. Strikt unterschieden werden
Retrodiktion, Bayes-Aktualisierung, semantische Rückbestimmung,
zeitumkehr\-symmetrische Dynamik, zweiseitige Randbedingungen, ontische
Rückwirkung, Rückwärts\-signalisierung und geschlossene zeitartige Kurven. Die
Unterscheidung wird interventionistisch formuliert, an avancierten und
retardierten Lösungen, Prä- und Postselektion, schwachen Werten,
Delayed-Choice- und Quantum-Eraser-Experimenten, Bell-Korrelationen,
thermodynamischem Zeitpfeil und relativistischer Kausalstruktur geprüft. Ein
Satz zeigt, dass das hier definierte autonome Mandelbrot-/QIK--VRT-Modell
keinen rückwärts gerichteten Zustandskanal enthält. Eine physikalische
Retrokausalität bleibt eine präzise formulierbare, aber nicht durch diese
Mathematik allein bewiesene Zusatzhypothese. Teil V bündelt daraus
Falsifikationskriterien, Ontologie, öffentlich fixierte Vorarbeiten,
Forschungsprogramm und ethische Konsequenzen.
\end{abstract}

\begin{otherlanguage}{english}
\begin{abstract}
Part I proves that the restricted Mandelbrot class \(\M_U=\M\cap U\) and its relative complement
\(\K_U=U\setminus\M_U\) exhaust exactly a previously fixed domain \(U\).
Calling \(U\) a \emph{model universe} is a
definition; identifying it with the physical universe requires an additional
correspondence map and empirical tests. This paper proves the set-theoretic
partition, the exact escape-time exhaustion of the complement, and finite
representational quantizability of every bounded subset of finite-dimensional
Euclidean space at any prescribed positive accuracy. It also proves, by the
continuous countermodel \([0,1]^4\), that representational quantizability does
not entail ontic discreteness or physical quantization of spacetime.
Mathematical access to a black-hole interior is distinguished from causal
observation of that interior; the Planck scale is distinguished from an
experimentally established minimum length. These distinctions provide the
scientifically strongest formulation of the proposed ontology of universal
difference and its interpretation as a universal principle of creation.

Part II turns the static partition into a typed process model. The extended
dynamics \(F(c,z)=(c,z^2+c)\) generates trajectories, while finite gates
distinguish certified passage, continued examination, and certified blocking.
The exact classifier is proved shift-invariant; finite certificates refine
monotonically; and an exact factorization criterion is given for gate-preserving
process maps. Later evidence may reclassify an earlier trace without changing
any earlier state, so semantic back-determination is not physical
retrocausation. The relevant fixed point belongs to the classifier, not in
general to an individual orbit.

Part III supplies a dimensionally consistent bridge. It distinguishes a
mathematical metric, a Lorentzian spacetime metric, numerical measurement and
physical unit; introduces explicit reference scales; checks the dimensions of
the Einstein equation, action, Planck quantities and horizon entropy; and
states the additional observables and tests required for a physical theory.

Part IV treats retrocausality as a distinct family of physical and
epistemological hypotheses. It separates retrodiction, Bayesian updating,
semantic back-determination, time-reversal symmetry, two-boundary models,
ontic backward influence, operational backward signalling and closed timelike
curves. These notions are tested against advanced and retarded solutions,
pre- and post-selection, weak values, delayed-choice and quantum-eraser
experiments, Bell correlations, the thermodynamic arrow and relativistic
causal order. A theorem establishes that the autonomous Mandelbrot/QIK--VRT
process defined here contains no backward state channel. Physical
retrocausality therefore remains a precise additional hypothesis rather than
a consequence proved by the present mathematics. Part V collects the
resulting falsification criteria, ontology, public evidence, research
programme and ethical consequences.
\end{abstract}
\end{otherlanguage}

\textbf{Schlagworte:} Mandelbrot-Menge; relatives Komplement; Escape-Time;
Modelluniversum; Anschlussordnung; Wirkungsgate; Raumzeit; Diskretisierung;
Quantisierbarkeit;\linebreak Dimensionsanalyse; Quantengravitation; Ereignishorizont;
Retrokausalität; Retrodiktion; No-Signalling; Ontologie des Unterschieds;
\qikvrt.

\tableofcontents
\newpage

\part{Mathematisch-theoretische Grundlage}

\section{Ergebnis vorweg: Was zweifelsfrei gilt}

Der Ausdruck \emph{zweifelsfrei} wird in dieser Arbeit nicht rhetorisch,
sondern technisch verwendet: Eine Aussage ist genau dann als bewiesen
ausgewiesen, wenn sie aus angegebenen Definitionen und Voraussetzungen durch
einen nachvollziehbaren Schluss folgt. Die weitergehende physikalische oder
ontologische Bedeutung wird gesondert klassifiziert.

\statusbox{QGreen}{
\textbf{Mathematisch bewiesen.}
Für jeden festgelegten Grundbereich \(U\subseteq\C\) bilden
\(\M\cap U\) und sein relatives Komplement eine disjunkte und vollständige
Zerlegung von \(U\). Das Komplement ist die wachsende Vereinigung aller
endlichen Escape-Time-Mengen. Jeder beschränkte Bereich in \(\R^d\) ist bei
jeder Genauigkeit \(\eps>0\) endlich diskret darstellbar. Die
Komplementzerlegung ist dimensionsunabhängig; ihre Erhaltung unter
Abbildungen verlangt jedoch Unterscheidungstreue. Die exakte
Trajektorienklassifikation ist shift-invariant; endliche PASS- und
BLOCK-Zertifikate verfeinern sich monoton. Die angegebenen
Dimensionsgleichungen der etablierten Physik sind homogen. Retrodiktion,
semantische Rückbestimmung, ontische Retrokausalität und operative
Rückwärtssignalisierung sind definitorisch und mathematisch voneinander
trennbar. Das in diesem Dokument definierte autonome Prozessmodell enthält
keinen rückwärts gerichteten Zustandskanal.
}

\statusbox{QAmber}{
\textbf{Definiert beziehungsweise bedingt bewiesen.}
Wird der Grundbereich \(U\) als Modelluniversum definiert, dann ergeben
die eingeschränkte Mandelbrot-Klasse \(\M_U\) und ihre Komplementklasse
\(\K_U\) dieses Modelluniversum. Wird eine totale
Modellabbildung \(\Phi:X\to U\) für einen Raumzeitbereich \(X\) festgelegt,
so zerlegen die Urbilder der beiden Klassen auch \(X\) vollständig. Ob
\(\Phi\) die wirkliche Raumzeit physikalisch richtig beschreibt, ist eine
zusätzliche, empirisch offene Frage. Eine dynamische Übertragung ist nur dann
gate-erhaltend, wenn sie Statusunterschiede auf ihren Fasern nicht
zusammenführt. Eine physikalische Übertragung verlangt darüber hinaus
Referenzskalen, Invarianten, Dynamik, Observablen und Tests.
Eine retrokausale physikalische Erweiterung ist als Modell konstruierbar,
erfordert aber lokalisierte Variablen, eine interventionistisch definierte
spätere Einstellung, Lorentz- und Erhaltungskonsistenz, eine No-Signalling-
oder Kanalkapazitätsaussage und eine von Standardquantenmechanik
unterscheidbare Vorhersage.
}

\statusbox{QRed}{
\textbf{Nicht dadurch bewiesen.}
Nicht bewiesen sind die ontische Quantisierung der Raumzeit, eine kleinste
physikalische Länge, die Identität des komplexen Parameterraums mit dem
empirischen Kosmos, eine bestimmte abgeschlossene Theorie der
Quantengravitation, physikalische Retrokausalität aus späterer
Reklassifikation, ein Nash-Gleichgewicht ohne Spielstruktur oder das
universale Schöpfungsprinzip als experimentell bestätigtes Naturgesetz.
Ebenso wenig beweisen Zeitumkehrsymmetrie, Postselektion,
Delayed-Choice-Korrelationen, Bell-Verletzungen, avancierte mathematische
Lösungen oder die Existenz relativistischer Raumzeiten mit geschlossenen
zeitartigen Kurven für sich eine steuerbare Wirkung in die Vergangenheit.
}

Die kürzeste wissenschaftlich vollständige Fassung lautet daher:
\[
\boxed{
\begin{gathered}
\M_U\dotcup\K_U=U,\\[0.3em]
\widehat\chi\circ\sigma=\widehat\chi,\\[0.3em]
\text{darstellungsbezogene Quantisierbarkeit}
\neq
\text{ontische Quantisierung},\\[0.3em]
\text{Dimensionsverträglichkeit}
\neq
\text{physikalische Herleitung},\\[0.3em]
\text{mathematische Darstellbarkeit}
\neq
\text{kausale Beobachtbarkeit},\\[0.3em]
\text{semantische Rückbestimmung}
\neq
\text{physikalische Retrokausalität}
\neq
\text{Rückwärtssignalisierung}.
\end{gathered}}
\]

\section{Methodik und Geltungsebenen}

Die Hauptgefahr einer interdisziplinären Gesamtaussage ist nicht mangelnde
Reichweite, sondern die unmarkierte Verschiebung zwischen unterschiedlichen
Arten von Geltung. Die folgende Ordnung ist deshalb selbst eine Anwendung der
Ontologie des Unterschieds.

\begin{longtable}{>{\raggedright\arraybackslash}p{0.25\textwidth}
                  >{\raggedright\arraybackslash}p{0.32\textwidth}
                  >{\raggedright\arraybackslash}p{0.33\textwidth}}
\toprule
\textbf{Ebene} & \textbf{Funktion} & \textbf{Prüfkriterium} \\
\midrule
\endfirsthead
\toprule
\textbf{Ebene} & \textbf{Funktion} & \textbf{Prüfkriterium} \\
\midrule
\endhead
Mathematischer Satz & Aussage innerhalb expliziter Definitionen und Axiome & formaler Beweis oder Gegenbeispiel \\
Modellkonvention & Festlegung dessen, was im Modell ``Universum'', ``Komplement'' oder ``Quantisierer'' heißt & Eindeutigkeit und Widerspruchsfreiheit der Definition \\
Korrespondenz\-hypothese & Zuordnung zwischen Modellobjekten und physikalischen Ereignissen & Observablen, unterscheidende Vorhersagen, Messung \\
Empirische Aussage & Bericht über Messdaten oder reproduzierte Beobachtung & Unsicherheit, Kalibrierung, unabhängige Reproduktion \\
Interpretation & Ontologische Lesart empirisch äquivalenter Formalismen & begriffliche Konsistenz; für empirische Auszeichnung zusätzlich eine unterscheidende Vorhersage \\
Kausalaussage & Behauptung, dass eine Intervention eine andere Variable beeinflusst & Interventionsverteilung, Ausschluss gemeinsamer Ursachen, Zeit- und Ortszuordnung \\
Ontologische Deutung & Aussage über Bestimmtheit, Unterschied, Einheit und Wirklichkeit & begriffliche Tragfähigkeit, Selbstanwendung, Gegenargumente \\
Normative Folgerung & Aussage darüber, wie Technik und Forschung eingesetzt werden sollen & offengelegte Wertprämissen und Folgenabwägung \\
\bottomrule
\end{longtable}

Ein Zenodo-Archiv, ein Git-Commit und ein Hash belegen, welches Artefakt zu
welchem Zeitpunkt und in welcher Bytefolge öffentlich fixiert wurde. Sie
belegen nicht automatisch Peer Review, wissenschaftlichen Konsens oder die
physikalische Wahrheit jeder archivierten Behauptung. Diese Unterscheidung
wird in \cref{sec:evidenz} für die Vorarbeiten von Ingolf Lohmann explizit
angewandt.

\section{Die Mandelbrot-Menge und ihr Komplement}
\label{sec:mandelbrot}

\subsection{Definitionen}

Die Mandelbrot-Menge geht auf die Untersuchung komplexer quadratischer
Iterationen zurück und wurde durch Mandelbrot in ihrer fraktalen Bedeutung
sichtbar gemacht \citep{mandelbrot1980}.

\begin{definition}[Quadratische Iteration]
Für \(c\in\C\) sei
\[
  z_0(c)=0,
  \qquad
  z_{n+1}(c)=z_n(c)^2+c
  \quad(n\in\N_0).
\]
Die Mandelbrot-Menge ist
\[
  \M
  =
  \left\{
    c\in\C:
    \bigl(z_n(c)\bigr)_{n\in\N_0}
    \text{ ist beschränkt}
  \right\}.
\]
\end{definition}

\begin{definition}[Relatives Komplement]
Sei \(U\subseteq\C\) der ausdrücklich gewählte Grundbereich. Wir setzen
\[
  \M_U:=\M\cap U,
  \qquad
  \K_U:=U\setminus\M_U.
\]
\(\K_U\) heißt in diesem Dokument die zu \(U\) relative
Komplementärmenge der Mandelbrot-Menge.
\end{definition}

Ein Komplement ist ohne Grundbereich nicht eindeutig. Für \(U=\C\) lautet es
\(\C\setminus\M\); für ein endliches Bildfenster \(U\) dagegen nur
\(U\setminus\M_U\). Diese Relativität ist keine Schwäche, sondern die
Voraussetzung einer bestimmten Aussage.

\subsection{Der Komplementsatz}

\begin{satz}[Vollständige komplementäre Zerlegung]
\label{satz:komplement}
Für jedes \(U\subseteq\C\) gilt
\[
  U=\M_U\dotcup\K_U,
\]
also
\[
  \M_U\cup\K_U=U
  \qquad\text{und}\qquad
  \M_U\cap\K_U=\varnothing.
\]
Insbesondere gilt
\[
  \M\cup(\C\setminus\M)=\C.
\]
\end{satz}

\begin{proof}
Sei \(u\in U\). In klassischer Logik gilt entweder
\(u\in\M_U\) oder \(u\notin\M_U\). Im zweiten Fall ist nach Definition
\(u\in U\setminus\M_U=\K_U\). Also liegt jedes Element von \(U\) in der
Vereinigung. Umgekehrt sind beide Teilmengen in \(U\) enthalten. Damit ist die
Vereinigung gleich \(U\). Ein Element von \(\K_U\) ist definitionsgemäß kein
Element von \(\M_U\); der Schnitt ist daher leer.
\end{proof}

Der Satz beweist eine vollständige Partition. Er beweist nicht, dass die
Mandelbrot-Menge erst durch eine Visualisierung ihr Komplement erhält. Das
Komplement ist bereits durch \(U\) und \(\M_U\) eindeutig bestimmt.

\subsection{Escape-Time und konstruktiv zertifiziertes Außen}

\begin{definition}[Fluchtzeit]
Die Fluchtzeit eines Parameters \(c\in\C\) ist
\[
  \tau(c)
  :=
  \inf\{n\in\N_0:|z_n(c)|>2\}
  \in\N_0\cup\{\infty\},
\]
wobei das Infimum der leeren Menge als \(\infty\) definiert wird.
\end{definition}

\begin{lemma}[Fluchtradius]
\label{lemma:escape}
Sobald \(|z_n(c)|>2\) für ein \(n\) gilt, ist die Folge
\((z_k(c))_{k\in\N_0}\) unbeschränkt.
\end{lemma}

\begin{proof}
Ist \(|c|>2\), so gilt bereits \(|z_1|=|c|>2\). Für
\(r_k:=|z_k|\ge |c|>2\) folgt
\[
  r_{k+1}\ge r_k^2-|c|\ge r_k^2-r_k=r_k(r_k-1)>r_k.
\]
Ist dagegen \(|c|\le 2\) und \(r_n>2\), so gilt
\[
  r_{k+1}\ge r_k^2-2>r_k
  \qquad(k\ge n),
\]
denn \(r_k^2-r_k-2=(r_k-2)(r_k+1)>0\). Eine angenommene
endliche Grenze \(L>2\) dieser monoton wachsenden Folge wäre mit
\(L\ge L^2-2\) unvereinbar. Also ist die Folge unbeschränkt.
\end{proof}

\begin{korollar}
\label{kor:escape}
Es gilt exakt
\[
  \M=\{c\in\C:\tau(c)=\infty\},
  \qquad
  \C\setminus\M=\{c\in\C:\tau(c)<\infty\}.
\]
\end{korollar}

\begin{definition}[Endliche Escape-Mengen]
Für \(N\in\N_0\) sei
\[
  E_N(U):=\{c\in U:\tau(c)\le N\}.
\]
\end{definition}

\begin{satz}[Stufenweise Rekonstruktion des Komplements]
\label{satz:escape-union}
Die Folge \((E_N(U))_{N\in\N_0}\) ist monoton wachsend und
\[
  \K_U=\bigcup_{N=0}^{\infty}E_N(U).
\]
\end{satz}

\begin{proof}
Aus \(\tau(c)\le N\) folgt \(\tau(c)\le N+1\), also
\(E_N(U)\subseteq E_{N+1}(U)\). Nach \cref{kor:escape} liegt ein Punkt genau
dann im Komplement, wenn seine Fluchtzeit endlich ist. Jede endliche
Fluchtzeit ist kleiner oder gleich einem geeigneten \(N\); umgekehrt ist jeder
Punkt eines \(E_N(U)\) nach \cref{lemma:escape} außerhalb von \(\M\).
\end{proof}

\begin{bemerkung}[Was ein endliches Bild tatsächlich beweist]
Für ein endliches Raster \(G\subset U\) und eine endliche Iterationsgrenze
\(N\) sind die Punkte in \(E_N(U)\cap G\) als Außenpunkte zertifiziert. Für
die übrigen Rasterpunkte gilt jedoch nur
\[
G\setminus E_N(U)
=
(\M\cap G)
\cup
\{c\in G\setminus\M:\tau(c)>N\}.
\]
``Bis \(N\) nicht geflohen'' ist deshalb kein allgemeiner endlicher Beweis
für \(c\in\M\). Eine Escape-Time-Grafik visualisiert eine kontrollierte
Approximation. Visualisierung und exakte Mengenzugehörigkeit sind zu
unterscheiden; Fragen der Berechenbarkeit des Mandelbrot-Randes sind
zusätzlich subtil \citep{hertling2005}.
\end{bemerkung}

\section{Vom Grundbereich zum raumzeitlichen Modelluniversum}
\label{sec:modeluniversum}

\subsection{``Universum'' als definierter Grundbereich}

In der Mengenlehre kann \emph{Universum} den Grundbereich einer Betrachtung
bezeichnen. In der Kosmologie meint \emph{Universum} die physische Gesamtheit
von Raumzeit, Materie, Feldern und ihren Beziehungen. Beide Bedeutungen dürfen
nicht stillschweigend identifiziert werden.

\begin{definition}[Mandelbrot-Modelluniversum]
Ein ausdrücklich gewählter Grundbereich \(U\subseteq\C\) heißt in diesem
Dokument \emph{Mandelbrot-Modelluniversum}, wenn seine beiden Klassen
\(\M_U\) und \(\K_U\) als komplementäre Modellzustände gelesen werden.
\end{definition}

Aus \cref{satz:komplement} folgt unmittelbar:
\[
  \boxed{\text{Mandelbrot-Klasse }\M_U+
  \text{Komplementklasse }\K_U
  =\text{gesamtes Modelluniversum}.}
\]
Das Pluszeichen steht hier für disjunkte Vereinigung, nicht für numerische
Addition. Die Gleichung ist exakt. Ihre physikalische Interpretation ist eine
gesonderte Modellentscheidung.

\subsection{Raumzeitlich induzierte Komplementärmenge}

\begin{definition}[Korrespondenzabbildung]
Sei \(X\) die Menge der im Modell betrachteten Raumzeitereignisse und
\[
  \Phi:X\longrightarrow U
\]
eine totale Abbildung. Definiert werden
\[
  X_{\M}:=\Phi^{-1}(\M_U),
  \qquad
  X_{\K}:=\Phi^{-1}(\K_U).
\]
\(X_{\K}\) heißt die durch \(\Phi\) induzierte raumzeitliche
Komplementärmenge.
\end{definition}

\begin{satz}[Raumzeitliche Urbildpartition]
\label{satz:urbild}
Für jede totale Abbildung \(\Phi:X\to U\) gilt
\[
  X=X_{\M}\dotcup X_{\K}.
\]
\end{satz}

\begin{proof}
Urbilder erhalten Komplemente:
\[
  \Phi^{-1}(U\setminus\M_U)
  =X\setminus\Phi^{-1}(\M_U).
\]
Mit \cref{satz:komplement} folgt die behauptete disjunkte Vereinigung.
\end{proof}

Damit ist die folgende Formulierung streng zulässig:

\begin{quote}
Relativ zu einer ausdrücklich definierten totalen Modellabbildung ergeben die
Mandelbrot-Klasse und ihre raumzeitlich induzierte Komplementklasse den
vollständigen betrachteten Raumzeitbereich, also das Universum des Modells.
\end{quote}

Nicht zulässig wäre der unbegründete Kurzschluss, \(\C\), \(\M\) oder
\(\K_U\) seien ohne eine solche Abbildung bereits physische Raumzeit. Für eine
naturwissenschaftliche Identifikation muss \(\Phi\) mindestens folgende
Angaben besitzen:

\begin{enumerate}
  \item eine koordinaten- oder invariant definierte physikalische Domäne,
  \item die Zuordnung der Modellparameter zu messbaren Größen,
  \item eine Dynamik oder Übergangsregel,
  \item mindestens eine von etablierten Modellen unterscheidbare Vorhersage,
  \item ein Messverfahren samt Unsicherheitsmodell,
  \item unabhängige Reproduktion.
\end{enumerate}

\section{Quantisierbarkeit ist nicht Quantisierung}
\label{sec:quantisierbarkeit}

Das Wort \emph{quantisieren} hat mindestens zwei verschiedene fachliche
Bedeutungen. Numerik und Informationstechnik quantisieren Werte durch endliche
Raster oder Klassen. Physik quantisiert ein klassisches System durch eine
Zuordnung zu Zuständen, Operatoren und Observablen einer Quantentheorie. Die
zweite Bedeutung ist wesentlich stärker und nicht für jede klassische
Observable ohne Obstruktionen verfügbar \citep{groenewold1946}.

\subsection{Darstellungsbezogene Definition}

\begin{definition}[Endliche \(\eps\)-Quantisierbarkeit]
Eine Teilmenge \(K\) eines metrischen Raums \((X,d)\) heißt bei Genauigkeit
\(\eps>0\) \emph{endlich darstellungsbezogen quantisierbar}, wenn eine
endliche Menge \(G_{\eps}\subseteq X\) und eine Abbildung
\[
  Q_{\eps}:K\longrightarrow G_{\eps}
\]
existieren, so dass
\[
  d\bigl(x,Q_{\eps}(x)\bigr)\le\eps
  \qquad\text{für alle }x\in K.
\]
\end{definition}

\begin{satz}[Quantisierbarkeit beschränkter euklidischer Bereiche]
\label{satz:quantisierbar}
Jede beschränkte Teilmenge \(K\subset\R^d\) ist für jedes
\(\eps>0\) endlich darstellungsbezogen quantisierbar.
\end{satz}

\begin{proof}
Da \(K\) beschränkt ist, liegt \(K\) in einem Würfel \([-R,R]^d\).
Wähle \(h\le 2\eps/\sqrt d\). In jeder Koordinate existiert ein endliches
\(h\)-Netz des Intervalls \([-R,R]\), das beide Randpunkte enthält und
dessen aufeinanderfolgende Punkte höchstens Abstand \(h\) besitzen. Das
kartesische Produkt \(G_{\eps}\) dieser Netze ist endlich. Ordne jedem
\(x\in K\) einen nächstgelegenen Punkt \(Q_{\eps}(x)\in G_{\eps}\) zu.
In jeder Koordinate beträgt der Fehler höchstens \(h/2\), also
\[
  \|x-Q_{\eps}(x)\|_2
  \le \sqrt d\,\frac h2
  \le\eps.
\]
\end{proof}

\begin{korollar}[Endlichdimensionale normierte Räume]
Jede beschränkte Teilmenge eines endlichdimensionalen normierten Vektorraums
ist bei jeder Genauigkeit \(\eps>0\) endlich darstellungsbezogen
quantisierbar.
\end{korollar}

\begin{proof}
Nach Wahl einer Basis ist der Raum linear isomorph zu \(\R^d\). In endlicher
Dimension sind alle Normen äquivalent. Daher wird eine beschränkte Menge auf
eine euklidisch beschränkte Menge abgebildet, und das Raster aus
\cref{satz:quantisierbar} kann unter Berücksichtigung der
Normäquivalenzkonstanten zurückübertragen werden.
\end{proof}

In allgemeinen metrischen Räumen ist die Existenz einer solchen endlichen
\(\eps\)-Darstellung für jedes \(\eps>0\) äquivalent zur
Totalbeschränktheit. Global unbeschränkte oder beliebige
unendlichdimensionale Räume erfüllen diese Voraussetzung nicht automatisch.
Der Ausdruck \emph{stets quantisierbar} bedeutet in diesem Dokument deshalb
präzise:

\begin{quote}
Jeder beschränkte Ausschnitt eines endlichdimensionalen normierten Modell-
oder Beobachtungsraums ist bei jeder vorgegebenen positiven Genauigkeit
endlich diskret darstellbar.
\end{quote}

Für jede Teilmenge \(A\subseteq U\) existiert ferner in klassischer
Mengenlehre die exakte binäre Charakteristik
\[
  \chi_A:U\to\{0,1\},
  \qquad
  \chi_A(x)=
  \begin{cases}
  1,&x\in A,\\
  0,&x\notin A.
  \end{cases}
\]
Das ist eine logische Zwei-Klassen-Darstellung. Ihre mengentheoretische
Existenz garantiert nicht, dass die Mitgliedschaft für jede gegebene
Beschreibung algorithmisch in endlicher Zeit entscheidbar ist.

\subsection{Der formale Nichtschluss}

\begin{proposition}[Quantisierbarkeit impliziert keine ontische Diskretheit]
\label{prop:nichtschluss}
Aus der Existenz endlicher \(\eps\)-Quantisierer für alle
\(\eps>0\) folgt weder eine kleinste positive Distanz noch die ontische
Diskretheit des dargestellten Raums.
\end{proposition}

\begin{proof}
Der metrische Raum \(K=[0,1]^4\) mit der euklidischen Metrik ist kontinuierlich
und besitzt keine kleinste positive Distanz: Zu jeder positiven Distanz
zwischen zwei Punkten existieren in \(K\) noch kleinere positive Distanzen.
Dennoch ist \(K\) nach \cref{satz:quantisierbar} bei jeder Genauigkeit
\(\eps>0\) endlich darstellungsbezogen quantisierbar. Damit ist \(K\) ein
Gegenmodell zur behaupteten Implikation.
\end{proof}

\statusbox{QLight}{
\textbf{Der entscheidende Unterschied.}
Ein Raster ist eine Eigenschaft der Darstellung und ihrer Auflösung. Ein
diskretes Spektrum eines physikalischen Operators wäre eine Eigenschaft einer
bestimmten Theorie. Eine fundamentale Raumzeitkörnung wäre eine Aussage über
die Wirklichkeit. Keine dieser drei Ebenen darf ohne Brückenbeweis in die
nächste überführt werden.
}

Regge zeigte, wie gekrümmte Geometrie durch Simplizialkomplexe approximiert
werden kann \citep{regge1961}. Dies belegt die Leistungsfähigkeit diskreter
Darstellungen klassischer Geometrie, nicht die physische Existenz elementarer
Raumzeitzellen. Umgekehrt zeigen Ansätze von Snyder und der Causal-Set-Theorie,
dass fundamental diskrete oder lokal endliche Raumzeitmodelle mathematisch
denkbar sind \citep{snyder1947,bombelli1987}. Ihre Denkbarkeit ist jedoch
nicht mit empirischer Entscheidung gleichzusetzen.

\section{Geltung in beliebiger Dimension und unter Vektorisierung}
\label{sec:dimension}

\begin{satz}[Dimensionsunabhängigkeit]
Für jede Menge \(X\), unabhängig von Dimension, Topologie, Metrik oder
Vektorraumstruktur, und jede Teilmenge \(A\subseteq X\) gilt
\[
  X=A\dotcup(X\setminus A).
\]
\end{satz}

Der Satz benutzt nur Zugehörigkeit und Nichtzugehörigkeit. In diesem präzisen
Sinn ist die Komplementstruktur in jeder endlichen oder unendlichen Dimension
anwendbar. Eine Aussage über \emph{jede Vektorisierung} verlangt dagegen eine
Zusatzbedingung.

\begin{satz}[Komplementerhaltung unter Abbildungen]
\label{satz:abbildung}
Seien \(A\subseteq X\) und \(f:X\to Y\). Dann gilt stets
\[
  f(X)\setminus f(A)\subseteq f(X\setminus A).
\]
Die Gleichheit
\[
  f(X\setminus A)=f(X)\setminus f(A)
\]
gilt genau dann, wenn
\[
  f(A)\cap f(X\setminus A)=\varnothing.
\]
Injektivität von \(f\) ist hinreichend. Ist \(f:X\to Y\) bijektiv, so gilt
\[
  f(X\setminus A)=Y\setminus f(A).
\]
\end{satz}

\begin{proof}
Ist \(y\in f(X)\setminus f(A)\), besitzt \(y\) ein Urbild in \(X\), aber
keines in \(A\); also liegt ein Urbild in \(X\setminus A\). Dies beweist die
Inklusion. Die umgekehrte Inklusion gilt genau dann, wenn kein Bildpunkt
zugleich von \(A\) und von \(X\setminus A\) getroffen wird. Injektivität
schließt ein solches Zusammenfallen aus. Bei Bijektivität ist zusätzlich
\(f(X)=Y\).
\end{proof}

Eine nichtinjektive Projektion, Dimensionsreduktion oder Feature-Vektorisierung
kann also Mandelbrot- und Komplementpunkte auf denselben Vektor abbilden. Die
universale Aussage lautet nicht ``jede Darstellung bewahrt jeden
Unterschied'', sondern:

\begin{quote}
Die binäre Komplementpartition ist von einer vorhandenen Dimensionsstruktur
unabhängig. Ihre eindeutige Übertragung unter einer Vektorisierung wird nur
durch injektive beziehungsweise anderweitig klassentrennende Darstellungen
garantiert. Urbilder einer totalen Modellabbildung bewahren sie stets.
\end{quote}

\part{Rekursive Anschlussordnung und selbstbegrenzte Klassifikation}

\section{Vom statischen Komplement zum geprüften Prozess}
\label{sec:dynamische-synthese}

Der Komplementsatz
\[
  U=\M_U\dotcup\K_U
\]
beschreibt zunächst eine statische Partition des Parameterraums. Die
quadratische Iteration enthält zusätzlich eine dynamische Struktur: Zu jedem
Parameter \(c\) gehört eine vollständige Anschlussgeschichte. Erst auf dieser
Ebene werden endliche Prüfung, fortgesetzte Unbestimmtheit, zertifizierter
Ausschluss und asymptotische Klassifikation unterscheidbar.

\subsection{Erweiterter Zustandsraum und Mandelbrot-Trajektorien}

\begin{definition}[Erweiterte Mandelbrot-Dynamik]
Sei \(U\subseteq\C\). Auf dem erweiterten Zustandsraum
\[
  Y_U:=U\times\C
\]
definieren wir
\[
  F:Y_U\longrightarrow Y_U,
  \qquad
  F(c,z):=(c,z^2+c),
\]
sowie die Initialeinbettung
\[
  \iota:U\longrightarrow Y_U,
  \qquad
  \iota(c):=(c,0).
\]
Die zu \(c\) gehörende Zustandstrajektorie ist
\[
  \Gamma_c:=\bigl(F^n(\iota(c))\bigr)_{n\in\N_0}.
\]
Für die Projektion \(\pi_z(c,z):=z\) gilt
\[
  \pi_z\!\left(F^n(\iota(c))\right)=z_n(c).
\]
\end{definition}

Diese Formulierung trennt den konstanten Parameter \(c\) vom iterierten
Zustand \(z_n\). Die Rekursion wird damit zu einem autonomen diskreten
dynamischen System.

\begin{definition}[Trajektorienraum und Shift]
Sei
\[
  \Sigma_U:=\left\{\sigma^k\Gamma_c:c\in U,\ k\in\N_0\right\},
\]
wobei der Linksshift durch
\[
  \sigma(x_0,x_1,x_2,\ldots):=(x_1,x_2,x_3,\ldots)
\]
gegeben ist.
\end{definition}

\begin{definition}[Exakte Trajektorienklassifikation]
Für \(\gamma=(x_n)_{n\in\N_0}\in\Sigma_U\) setzen wir
\[
  \widehat\chi(\gamma):=
  \begin{cases}
    \PASS,&\sup_{n\in\N_0}|\pi_z(x_n)|<\infty,\\[0.25em]
    \BLOCK,&\sup_{n\in\N_0}|\pi_z(x_n)|=\infty.
  \end{cases}
\]
Für unverschobene Trajektorien gilt somit
\[
  \widehat\chi(\Gamma_c)=
  \begin{cases}
    \PASS,&c\in\M_U,\\
    \BLOCK,&c\in\K_U.
  \end{cases}
\]
\end{definition}

\begin{satz}[Shift-Invarianz der exakten Klassifikation]
\label{satz:shift-invarianz}
Für alle \(\gamma\in\Sigma_U\) und \(k\in\N_0\) gilt
\[
  \widehat\chi(\sigma^k\gamma)=\widehat\chi(\gamma).
\]
\end{satz}

\begin{proof}
Das Entfernen endlich vieler Folgenglieder ändert weder Beschränktheit noch
Unbeschränktheit einer Folge. Daher besitzen \(\gamma\) und jeder endliche
Tail \(\sigma^k\gamma\) denselben Klassifikationsstatus.
\end{proof}

Sei \(\mathcal S\) der Pullback-Operator auf statuswertigen Funktionen,
\[
  (\mathcal S h)(\gamma):=h(\sigma\gamma).
\]
Aus \cref{satz:shift-invarianz} folgt
\[
  \boxed{\mathcal S\widehat\chi=\widehat\chi},
  \qquad
  \widehat\chi\in\Fix(\mathcal S).
\]
Die exakte Klassifikation ist damit ein Fixpunkt des auf
\emph{Klassifikatoren} wirkenden Shiftoperators.

\begin{bemerkung}[Kein allgemeiner Orbit-Fixpunkt]
Diese Fixpunktaussage betrifft nicht die einzelnen Zustandstrajektorien. Aus
\(\sigma\Gamma_c=\Gamma_c\) würde wegen \(z_0=0\) bereits \(z_1=c=0\)
folgen. Unter den auf \(z_0=0\) initialisierten Mandelbrot-Trajektorien ist
daher nur \(c=0\) shift-fixiert. Andere Parameter können periodische,
präperiodische, beschränkte nichtperiodische oder unbeschränkte Verläufe
erzeugen. Der mathematisch tragfähige Fixpunkt ist die unter endlicher
Verschiebung stabile Klassifikationsordnung, nicht ein vermeintlich
konvergierender Einzelorbit.
\end{bemerkung}

\subsection{Endliche Gate-Zustände}

Ein endlicher Rechenlauf kann Flucht positiv zertifizieren, im Allgemeinen
aber nicht allein aus ausgebliebener Flucht die Zugehörigkeit zur
Mandelbrot-Menge beweisen. Für eine wissenschaftlich korrekte
Dreizustandslogik muss \(\PASS\) deshalb ebenfalls an ein positives
Zertifikat gebunden werden.

\begin{definition}[Zulässige PASS-Zertifikate]
Eine Familie \((P_N)_{N\in\N_0}\) heißt zulässige Familie endlicher
PASS-Zertifikate, wenn
\[
  P_N\subseteq\M_U
  \qquad\text{und}\qquad
  P_N\subseteq P_{N+1}
\]
für alle \(N\) gelten. Die Aufnahme \(c\in P_N\) muss durch einen gültigen
Mitgliedschaftsnachweis begründet sein und darf nicht allein auf
\(\tau(c)>N\) beruhen.
\end{definition}

\begin{definition}[Endliches Anschlussgate]
Für \(N\in\N_0\) sei
\[
  \mathcal A_N:U\longrightarrow\{\PASS,\CONTINUE,\BLOCK\}
\]
durch
\[
  \mathcal A_N(c):=
  \begin{cases}
    \BLOCK,&\tau(c)\le N,\\
    \PASS,&c\in P_N,\\
    \CONTINUE,&\text{sonst}
  \end{cases}
\]
definiert.
\end{definition}

Da \(P_N\subseteq\M_U\) und \(E_N(U)\subseteq\K_U\) gilt, können die
ersten beiden Fälle nicht gleichzeitig eintreten.

\begin{satz}[Korrektheit und monotone Verfeinerung]
\label{satz:gate-korrektheit}
Das endliche Anschlussgate besitzt folgende Eigenschaften:
\begin{enumerate}
  \item \(\mathcal A_N(c)=\PASS\) impliziert \(c\in\M_U\);
  \item \(\mathcal A_N(c)=\BLOCK\) impliziert \(c\in\K_U\);
  \item ein einmal erreichter Status \(\PASS\) oder \(\BLOCK\) bleibt bei
        allen späteren Stufen erhalten;
  \item \(\CONTINUE\) bezeichnet weder Mitgliedschaft noch
        Nichtmitgliedschaft, sondern den gegenwärtig nicht abgeschlossenen
        Prüfzustand.
\end{enumerate}
\end{satz}

\begin{proof}
Die ersten beiden Aussagen folgen aus \(P_N\subseteq\M_U\) und dem
Fluchtradiuslemma. Wegen \(P_N\subseteq P_{N+1}\) bleibt ein
PASS-Zertifikat gültig. Wegen \(E_N(U)\subseteq E_{N+1}(U)\) bleibt auch
ein BLOCK-Zertifikat gültig. Für einen Punkt außerhalb beider
Zertifikatsmengen ist auf Stufe \(N\) noch keine terminale Folgerung zulässig.
\end{proof}

Die Informationsordnung der Statuswerte ist daher nicht linear, sondern
verzweigt:
\[
  \CONTINUE\preccurlyeq\PASS,
  \qquad
  \CONTINUE\preccurlyeq\BLOCK,
\]
während \(\PASS\) und \(\BLOCK\) unvergleichbar bleiben.

\begin{satz}[Exterior-Vollständigkeit und bedingte Gesamtvollständigkeit]
\label{satz:gate-vollstaendigkeit}
Für jedes \(c\in\K_U\) existiert ein endliches \(N_0\), sodass
\[
  \mathcal A_N(c)=\BLOCK
  \qquad\text{für alle }N\ge N_0.
\]
Für jedes \(c\in\M_U\) wird dagegen genau dann nach endlich vielen Stufen
\(\PASS\) erreicht, wenn
\[
  c\in\bigcup_{N=0}^{\infty}P_N.
\]
Insbesondere wird das Gate auf ganz \(U\) genau dann schließlich terminal,
wenn \(\bigcup_NP_N=\M_U\).
\end{satz}

\begin{proof}
Für \(c\in\K_U\) ist \(\tau(c)<\infty\). Mit \(N_0=\tau(c)\) gilt daher
\(c\in E_N(U)\) für alle \(N\ge N_0\). Die PASS-Aussage folgt unmittelbar
aus der Definition der wachsenden Zertifikatsfamilie.
\end{proof}

Ohne einen zusätzlich bewiesenen Vollständigkeitssatz
\(\bigcup_NP_N=\M_U\) dürfen Punkte der Mandelbrot-Menge dauerhaft im Status
\(\CONTINUE\) verbleiben. Das ist kein Fehler des Gates, sondern die
korrekte Darstellung begrenzter endlicher Evidenz.

\subsection{Der formale Null-Eins-Übergang}

Die Entstehung eines endlichen Fluchtzertifikats kann als monotone boolesche
Folge dargestellt werden:
\[
  b_N(c):=\mathbf 1_{E_N(U)}(c),
  \qquad b_N(c)\le b_{N+1}(c).
\]
Für \(c\in\K_U\) erfolgt bei \(N=\tau(c)\) genau ein irreversibler
Zertifikatsübergang
\[
  0\longrightarrow1.
\]
Für \(c\in\M_U\) bleibt \(b_N(c)=0\) für alle \(N\).

\begin{bemerkung}
Vor der Flucht bedeutet \(b_N(c)=0\) lediglich: Bis Stufe \(N\) liegt kein
Fluchtzertifikat vor. Der Wert \(0\) ist deshalb kein positiver Beweis für
\(c\in\M_U\). Der Übergang \(0\to1\) bezeichnet hier präzise die Entstehung
eines unterscheidbaren Außenzertifikats und nicht ohne weitere Prämissen einen
kosmologischen Schöpfungsakt. Ebenso startet der konkrete Mandelbrot-Orbit bei
\(z_0=0\) mit \(z_1=c\), nicht allgemein mit \(z_1=1\); \(0\to1\) ist hier
ein Code für das Auftreten eines Unterschieds, kein Ersatz für die Dynamik.
\end{bemerkung}

\subsection{Grenze als Instabilitätsort der Klassifikation}

\begin{satz}[Grenze und Statusdiskontinuität]
\label{satz:grenzklassifikation}
Versehen wir \(\{\PASS,\BLOCK\}\) mit der diskreten Topologie. Dann ist
die exakte Parameterklassifikation
\[
  \chi_U(c):=
  \begin{cases}
    \PASS,&c\in\M_U,\\
    \BLOCK,&c\in\K_U
  \end{cases}
\]
genau an den Punkten des relativen Randes \(\partial_U\M_U\) unstetig.
\end{satz}

\begin{proof}
Im relativen Inneren von \(\M_U\) ist \(\chi_U\) lokal konstant \(\PASS\).
Da \(\M\) abgeschlossen ist, ist \(\K_U\) relativ offen; dort ist
\(\chi_U\) lokal konstant \(\BLOCK\). An jedem Randpunkt schneidet jede
relative Umgebung beide Klassen. Das Urbild eines einzelnen diskreten Status
kann dort keine Umgebung des Randpunktes enthalten. Folglich ist \(\chi_U\)
genau am Rand unstetig.
\end{proof}

Die fraktale Grenze ist damit ein mathematischer Entscheidungshorizont:
Beliebig kleine Parameteränderungen können dort unterschiedliche exakte
Klassen treffen. Der Ausdruck bezeichnet keine Gleichsetzung mit einem
physikalischen Ereignishorizont.

\subsection{Semantische Rückbestimmung ohne Retrokausalität}

Für einen Parameter \(c\) seien
\[
  \Pi_N(c):=(z_0(c),z_1(c),\ldots,z_N(c))
\]
die bis zur Stufe \(N\) bekannten Trajektorienpräfixe. Eine spätere Prüfstufe
ersetzt \(\Pi_N(c)\) nicht, sondern erweitert sie zu \(\Pi_{N+1}(c)\).

\begin{proposition}[Epistemische, nicht dynamische Rückbestimmung]
\label{prop:rueckbestimmung}
Wechselt der Gate-Status bei wachsendem \(N\) von \(\CONTINUE\) zu
\(\PASS\) oder \(\BLOCK\), so ändert sich kein früherer Zustand
\(z_k(c)\). Geändert wird ausschließlich die anhand einer erweiterten
Informationsbasis zulässige Klassifikation der bereits vorliegenden Spur.
\end{proposition}

\begin{proof}
Die Rekursion bestimmt jedes \(z_k(c)\) eindeutig aus \(c\) und \(z_0=0\).
Eine spätere Berechnung verändert diese Werte nicht. Sie kann lediglich ein
zuvor fehlendes Flucht- oder Mitgliedschaftszertifikat verfügbar machen.
\end{proof}

Die logische Asymmetrie lautet
\[
  c\notin\M\iff\exists n\in\N_0:\ |z_n(c)|>2,
  \qquad
  c\in\M\iff\forall n\in\N_0:\ |z_n(c)|\le2.
\]
Ein späteres Glied kann daher eine zuvor offene Außenfrage endlich
entscheiden. Dies ist eine \emph{semantische Rückbestimmung} der gesamten
Spur, keine physikalische Wirkung aus der Zukunft auf frühere Systemzustände.
Die vollständige begriffliche, quantenphysikalische und relativistische
Abgrenzung folgt in \cref{sec:retrokausalitaet}.

\subsection{Das Prozessmodelluniversum}

\begin{definition}[Mandelbrot-Prozessmodelluniversum]
Das dynamische Mandelbrot-Prozessmodelluniversum über \(U\) ist das Tupel
\[
  \mathfrak U_U:=
  \left(Y_U,F,\iota,\Sigma_U,\sigma,
  (\mathcal A_N)_{N\in\N_0},\widehat\chi,
  \M_U,\K_U,\partial_U\M_U\right).
\]
\end{definition}

Dieses Tupel enthält nicht nur die fertige Partition, sondern zugleich
\[
\begin{gathered}
\text{Initialisierung}
\longrightarrow\text{Iteration}
\longrightarrow\text{endliche Evidenz}
\\
\longrightarrow\text{Gate-Status}
\longrightarrow\text{exakte Klasse}
\longrightarrow\text{Grenzstruktur}.
\end{gathered}
\]
In diesem präzisen Sinn erzeugt dieselbe Rekursion die Anschlussgeschichte,
die zertifizierbare Außenklasse, die beschränkte Klasse und die gemeinsame
Grenze. Sie erzeugt jedoch nicht allein den gewählten Grundbereich \(U\), die
Interpretation der Klassen oder eine physikalische Raumzeitkorrespondenz.

\section{Vom Mandelbrot-Prozess zu einem allgemeinen Wirkraum}
\label{sec:wirkraum-korrespondenz}

\subsection{Typisierte Wirkraumstruktur}

Die symbolische Schreibweise
\[
  \mathcal W=\mathcal A
  (\mathcal R\times\mathcal T\times\mathcal C_{\mathrm{caus}}\times\mathcal I)
\]
wird erst dann mathematisch bestimmt, wenn die beteiligten Mengen und der Typ
von \(\mathcal A\) angegeben sind.

\begin{definition}[Allgemeines Wirkraummodell]
Seien \(\mathcal R\), \(\mathcal T\), \(\mathcal C_{\mathrm{caus}}\) und
\(\mathcal I\) Mengen geometrisch-relationaler, temporaler, kausaler und
informationeller Zustandsdaten. Der rohe Zustandsraum sei
\[
  S:=\mathcal R\times\mathcal T\times
  \mathcal C_{\mathrm{caus}}\times\mathcal I.
\]
Ein allgemeines Wirkraummodell ist ein Tupel
\[
  \mathsf G:=\bigl(S,G,\Sigma_G,(\mathcal B_N)_{N\in\N_0},
  \mathcal B_\infty\bigr),
\]
wobei \(G:S\to S\) eine Übergangsdynamik, \(\Sigma_G\) der von \(G\)
erzeugte Trajektorienraum, \(\mathcal B_N\) endliche Gates und
\(\mathcal B_\infty\) die terminale Klassifikation sind.
\end{definition}

Das kartesische Produkt erzeugt weder Kausalität noch Dynamik. Diese entstehen
erst durch Relationen auf den Faktoren, Kompatibilitätsbedingungen, die
Übergangsabbildung \(G\) und die Gate-Regeln.

\subsection{Semikonjugation und Gate-Erhaltung}

Sei
\[
  Y_U^{\mathrm{rch}}:=
  \{F^n(\iota(c)):c\in U,\ n\in\N_0\}
\]
der erreichbare Mandelbrot-Zustandsraum.

\begin{definition}[Dynamische Semikonjugation]
Eine Abbildung \(\phi:Y_U^{\mathrm{rch}}\to S\) heißt Semikonjugation von
\(F\) nach \(G\), wenn
\[
  \boxed{G\circ\phi=\phi\circ F}.
\]
Sie induziert auf Trajektorien die koordinatenweise Abbildung
\[
  \Phi_\infty:\Sigma_U\longrightarrow\Sigma_G,
  \qquad
  \Phi_\infty((x_n)_n):=(\phi(x_n))_n,
\]
und erfüllt \(\Phi_\infty\circ\sigma=\sigma\circ\Phi_\infty\).
\end{definition}

Eine Semikonjugation erhält die Übergangsfolge. Sie darf jedoch verschiedene
Mandelbrot-Zustände identifizieren und beweist daher weder metrische noch
physikalische Gleichwertigkeit.

\begin{definition}[Gate-erhaltende Prozessabbildung]
Die Semikonjugation \(\phi\) heißt gate-erhaltend, wenn für alle
\(\gamma\in\Sigma_U\) und \(N\in\N_0\)
\[
  \mathcal B_N(\Phi_\infty(\gamma))=\mathcal A_N(\gamma)
\]
gilt und zusätzlich
\[
  \mathcal B_\infty(\Phi_\infty(\gamma))=\widehat\chi(\gamma).
\]
Dabei bezeichnet \(\mathcal A_N(\gamma)\) das auf den in \(\gamma\)
gespeicherten Parameter angewandte Mandelbrot-Gate.
\end{definition}

\begin{satz}[Faktorisierungskriterium für Gate-Erhaltung]
\label{satz:gate-faktorisierung}
Sei \(\Phi_\infty:\Sigma_U\to\Sigma_G\) eine Trajektorienabbildung. Für
ein festes \(N\) existiert genau dann ein deterministisches Gate
\(\mathcal B_N\) auf \(\Phi_\infty(\Sigma_U)\) mit
\[
  \mathcal A_N=\mathcal B_N\circ\Phi_\infty,
\]
wenn \(\mathcal A_N\) auf jeder Faser von \(\Phi_\infty\) konstant ist.
\end{satz}

\begin{proof}
Existiert \(\mathcal B_N\), folgt die Faserkonstanz durch Einsetzen. Gilt sie
umgekehrt, definieren wir für \(\eta\in\Phi_\infty(\Sigma_U)\)
\(\mathcal B_N(\eta):=\mathcal A_N(\gamma)\), wobei \(\gamma\) ein
beliebiges Urbild von \(\eta\) ist. Die Faserkonstanz garantiert die
Unabhängigkeit von der Wahl des Urbildes.
\end{proof}

Identifiziert \(\Phi_\infty\) zwei Trajektorien mit unterschiedlichen
PASS-, CONTINUE- oder BLOCK-Zuständen, kann auf dem Bildraum kein
deterministisches Gate die ursprüngliche Klassifikation vollständig erhalten.
Das minimale Kriterium einer nicht bloß metaphorischen Übertragung lautet
daher
\[
  \boxed{\text{Dynamikerhaltung}+\text{Gate-Erhaltung}+
  \text{Klassentrennung}.}
\]
Für eine physikalische Korrespondenz treten Dimensions-, Metrik-, Symmetrie-,
Kausal- und Observablenerhaltung hinzu.

\subsection{Erreichbarkeitsordnung und physikalische Kausalität}

Auf einem diskreten dynamischen System kann die Erreichbarkeitsrelation
\[
  x\preccurlyeq_F y
  \quad:\Longleftrightarrow\quad
  \exists k\in\N_0:\ y=F^k(x)
\]
definiert werden. Aus der Semikonjugation folgt
\[
  x\preccurlyeq_F y
  \Longrightarrow
  \phi(x)\preccurlyeq_G\phi(y).
\]
Diese Relation ist zunächst eine rechnerische Anschlussordnung. Sie ist nicht
automatisch die physikalische Kausalrelation einer Lorentz-Mannigfaltigkeit.
Dafür müsste zusätzlich gezeigt werden, dass die Bildzustände
Raumzeitereignisse repräsentieren und zum Beispiel
\[
  x\preccurlyeq_F y
  \Longrightarrow
  \phi(y)\in J^+(\phi(x))
\]
mit einer empirisch bestimmten Lichtkegelstruktur gilt.

\subsection{Limes, Fixpunkt und Nash-Begriff}

\statusbox{QAmber}{
\textbf{Limes- und Fixpunktwarnung.}
Der Ausdruck
\[
  \Fix\!\left(\mathcal A\!\left(\lim_{n\to\infty}G^n(s)\right)\right)
\]
ist nur definiert, wenn der Limes in einer angegebenen Topologie existiert,
der Typ von \(\mathcal A\) feststeht und der Operator des Fixpunkts benannt
ist. Chaotische, periodische und bloß beschränkte Trajektorien müssen keinen
Zustandslimes besitzen. Eine stabile Klassifikation kann dennoch existieren.
}

Die für die Mandelbrot-Dynamik ohne Orbitkonvergenz gültige Verdichtung ist
\[
  \boxed{\widehat\chi\in\Fix(\mathcal S)},
  \qquad (\mathcal S h)(\gamma)=h(\sigma\gamma).
\]

Der Ausdruck \emph{Nash-stabil} verlangt darüber hinaus ein explizites Spiel
aus Spielern, Strategieräumen und Nutzenfunktionen sowie für eine
Stabilitätsaussage eine Anpassungsdynamik. Ein Strategieprofil \(s^*\) ist
Nash-Gleichgewicht genau dann, wenn kein Spieler durch einseitige Abweichung
seinen Nutzen erhöhen kann \citep{nash1950}. Aus Mandelbrot-Iteration,
Semikonjugation oder Shift-Fixpunkt allein folgt kein Nash-Gleichgewicht.
Ohne Spielstruktur ist daher \emph{shift-invariante Klassifikation} die hier
bewiesene und wissenschaftlich präzise Bezeichnung.

\subsection{Selbstanwendung als Wahrheits- und Wirkungsgate}

Die Prozessgrammatik lässt sich auf ihre eigenen Aussagen anwenden. Für eine
Behauptung \(p\), eine Evidenzmenge \(E\) und ein explizites Prüfverfahren
\(V\) sei
\[
  H(p,E,V)\in\{\PASS,\CONTINUE,\BLOCK\}.
\]
Ein soundes Gate erfüllt mindestens
\[
\begin{aligned}
H(p,E,V)=\PASS
&\Longrightarrow\text{gültiger Beleg für }p,\\
H(p,E,V)=\BLOCK
&\Longrightarrow\text{Widerlegung oder verletzte Freigabebedingung},\\
H(p,E,V)=\CONTINUE
&\Longrightarrow\text{vorhandene Evidenz ist nicht terminal}.
\end{aligned}
\]
Praktisch wird dazu die Aussage nach Gegenstand, Kontext, Quelle, Beleg,
Schluss, Grenze und möglicher Wirkung zerlegt. Herkunft, Widerspruch,
Anschluss und Geltungsgrad werden getrennt geprüft; erst dann wird eine
Wirkung freigegeben oder blockiert. Das Gate erzeugt keine Wahrheit. Es
verhindert den Schluss von fehlender Widerlegung auf erwiesene Wahrheit und
macht die Geltungsgrenze revisionsfähig.

\subsection{Ergebnis des zweiten Teils}

Der dynamische Kern lässt sich ohne Kategorienvermischung zusammenfassen:
\[
\boxed{
\begin{gathered}
  c\longmapsto\Gamma_c
  \longmapsto(\mathcal A_N(c))_{N\in\N_0}
  \longmapsto\widehat\chi(\Gamma_c),\\[0.3em]
  \widehat\chi\circ\sigma=\widehat\chi,\\[0.3em]
  G\circ\phi=\phi\circ F,\\[0.3em]
  \mathcal B_N\circ\Phi_\infty=\mathcal A_N
  \quad\text{nur bei gate-konstanten Fasern}.
\end{gathered}}
\]
Bewiesen sind damit die dynamische Prozessdarstellung, die shift-invariante
Endklassifikation, die sounde endliche Gate-Logik, die semantische
Rückbestimmung und die exakten Bedingungen einer gate-erhaltenden
Semikonjugation. Noch nicht bewiesen ist, dass ein bestimmter physikalischer
Wirkraum eine solche Semikonjugation mit den zusätzlich erforderlichen
metrischen, dimensionalen, kausalen und beobachtbaren Eigenschaften besitzt.

\part{Dimensionsrichtige physikalische Konkretisierung}

\section{Ereignishorizont, theoretischer Blick und Planck-Skala}
\label{sec:physik}

\subsection{Mathematischer Zugang hinter den Horizont}

In einer geeigneten konformen Vervollständigung einer asymptotisch flachen
Raumzeit wird der zukünftige Ereignishorizont unter den üblichen globalen
Voraussetzungen als Grenze des kausalen Vergangenheitsbereichs der
zukünftigen Nullunendlichkeit definiert \citep{wald1984}:
\[
  \mathcal H^+=\partial J^{-}(\mathcal I^+).
\]
In Schwarzschild-Koordinaten erscheinen die Metrikkomponenten bei
\(r=2GM/c^2\) singulär. Reguläre Eddington--Finkelstein- und
Kruskal--Szekeres-Koordinaten zeigen, dass es sich dort in der
Schwarzschild-Lösung um eine Koordinatensingularität handelt und die
klassische Lösung regulär über den Horizont fortgesetzt werden kann
\citep{finkelstein1958,kruskal1960}. Damit lässt sich das Innere eines
idealisierten klassischen Schwarzen Lochs mathematisch beschreiben und in
Kausaldiagrammen darstellen.

Diese mathematische Fortsetzbarkeit bestätigt die innere Konsistenz der
idealisierten klassischen Lösung. Sie ist kein empirischer Nachweis dafür,
dass das Innere eines realen astrophysikalischen Schwarzen Lochs in allen
Regimen exakt durch diese Fortsetzung beschrieben wird.

Dies rechtfertigt die Metapher, mit dem \emph{geistigen Auge} hinter den
Ereignishorizont zu blicken, wenn sie exakt als theoretische Modellierung
verstanden wird. Sie bedeutet nicht, dass ein äußerer Beobachter Licht oder
Messdaten aus dem Inneren eines zukünftigen Ereignishorizonts empfangen kann.
Oppenheimer und Snyder beschrieben bereits den gravitativen Kollaps und seine
kausale Abschließung im klassischen Modell \citep{oppenheimer1939}.

Auch die Bilder des Event Horizon Telescope zeigen nicht das Innere. Sie
rekonstruieren horizon-nahe Radioemission und eine zentrale
Helligkeitsdepression, die mit dem erwarteten Schatten eines Kerr-Schwarzen
Lochs vereinbar ist \citep{eht2019i,eht2019vi}. Damit gilt:
\[
  \boxed{
  \text{mathematische Darstellbarkeit des Inneren}
  \neq
  \text{kausale Beobachtbarkeit des Inneren}.}
\]

\subsection{Was die Planck-Skala besagt}

Die Planck-Länge ist
\[
  \ell_{\mathrm P}=\sqrt{\frac{\hbar G}{c^3}}
  \approx 1.616255\times10^{-35}\,\mathrm m,
\]
wobei die Naturkonstanten nach dem CODATA-Referenzrahmen zu verwenden sind
\citep{codata2022}. Sie kombiniert die Quantenskala \(\hbar\), die
Gravitationskonstante \(G\) und die Lichtgeschwindigkeit \(c\). Deshalb ist
sie eine natürliche charakteristische Skala für die Frage, wie
Quantentheorie und dynamische Raumzeit zusammenwirken.

Aus dieser Dimensionskombination folgt nicht, dass \(\ell_{\mathrm P}\) eine
experimentell bestätigte kleinste Länge, eine räumliche Zelle oder ein
Ereignishorizont ist. Verschiedene Quantengravitationsprogramme formulieren
unterschiedliche Begriffe von Diskretheit, Minimalität und Kontinuum
\citep{hossenfelder2013}. Die Allgemeine Relativitätstheorie bleibt bei
gewöhnlichen Energien als effektive Feldtheorie sinnvoll behandelbar
\citep{donoghue1994}; die hochenergetische fundamentale Vervollständigung ist
dadurch nicht ausgewählt.

Der epistemische Kontrast ist gleichwohl fruchtbar:

\begin{itemize}
  \item Hinter einem klassischen Ereignishorizont ist die Geometrie innerhalb
  einer Theorie mathematisch fortsetzbar, während der äußere Messzugang kausal
  begrenzt ist.
  \item Im Planck-Regime fehlt eine empirisch bestätigte gemeinsame Theorie
  von Quantendynamik und Gravitation. Die Schwierigkeit liegt nicht in einem
  bewiesenen ``innerhalb'' einer kleinsten Zelle, sondern in Theorie,
  Skalenreichweite und experimenteller Zugänglichkeit.
\end{itemize}

Diese beiden Grenzen sind nicht symmetrische physikalische Horizonte. Sie sind
verschiedene Arten von Grenze: kausale Grenze einer gegebenen Raumzeitlösung
einerseits, offene Theorie- und Messgrenze andererseits.

\section{Vom dimensionslosen Modell zur messbaren Physik}
\label{sec:dimensionsbruecke}

\statusbox{QLight}{
\textbf{Status dieses Teils.}
Die folgenden Dimensionsgleichungen, Definitionen der Raumzeitmetrik,
Einstein-Gleichungen, Planck-Einheiten und Beziehungen der
Schwarzloch-Thermodynamik gehören zur etablierten mathematischen Physik. Die
Übertragung der rekursiven Mandelbrot- und Anschlussordnung auf diese Größen
ist dagegen eine ausdrücklich definierte Korrespondenzhypothese.
Dimensionsrichtigkeit ist notwendig, aber nicht hinreichend für physikalische
Wahrheit.
}

\subsection{Größe, Zahlenwert, Einheit und Dimension}

Eine physikalische Größe \(q\) ist von ihrem Zahlenwert und ihrer Einheit zu
unterscheiden. Relativ zu einer gewählten Einheit \(u_q\) gilt
\[
  q=\{q\}_{u_q}\,u_q.
\]
Der Zahlenwert \(\{q\}_{u_q}\) ändert sich bei einem Einheitenwechsel; die
physikalische Größe soll davon unberührt bleiben. Das Internationale
Einheitensystem legt die einheitliche metrologische Darstellung fest
\citep{bipm2026si}. Die Dimension einer Größe kann als Produkt von
Basisdimensionen geschrieben werden:
\[
  \dim(q)=\mathsf L^{\alpha}\mathsf T^{\beta}\mathsf M^{\gamma}
  \boldsymbol\Theta^{\delta}\cdots,
\]
wobei \(\mathsf L,\mathsf T,\mathsf M,\boldsymbol\Theta\) Länge, Zeit,
Masse und thermodynamische Temperatur bezeichnen.

\begin{proposition}[Dimensionshomogenität]
Eine einheitenunabhängige physikalische Gleichung muss dimensionshomogen
sein: Addierbare Terme besitzen dieselbe Dimension, und Argumente von
Funktionen wie \(\exp\), \(\log\), \(\sin\) oder \(\cos\) müssen
dimensionslos sein.
\end{proposition}

\begin{proof}
Bei einer unabhängigen Skalierung der Basiseinheiten transformiert jeder Term
entsprechend seinem Dimensionsvektor. Verschiedene Dimensionsvektoren erhalten
verschiedene Skalierungsfaktoren und können daher nicht einheitenunabhängig
addiert oder gleichgesetzt werden. Ein nicht dimensionsloses
Funktionsargument würde seinen Zahlenwert allein durch den Einheitenwechsel
ändern und damit fälschlich den Funktionswert verändern. Dies ist der Kern der
Dimensionsanalyse \citep{buckingham1914}.
\end{proof}

Die Umkehrung gilt nicht. Dimensionshomogenität schließt einen Teil
unzulässiger Formeln aus, wählt aber nicht eindeutig das richtige Naturgesetz:
\[
  \boxed{
  \text{Dimensionsverträglichkeit}
  \neq\text{dynamische Herleitung}
  \neq\text{empirische Bestätigung}.}
\]

\begin{longtable}{>{\raggedright\arraybackslash}p{0.29\textwidth}
                  >{\raggedright\arraybackslash}p{0.23\textwidth}
                  >{\raggedright\arraybackslash}p{0.38\textwidth}}
\toprule
\textbf{Größe} & \textbf{SI-Einheit} & \textbf{Rolle} \\
\midrule
\endfirsthead
\toprule
\textbf{Größe} & \textbf{SI-Einheit} & \textbf{Rolle} \\
\midrule
\endhead
räumliche Koordinate \(x^i\) & \(\mathrm m\) & koordinatenabhängige Darstellung \\
Koordinate \(x^0=ct\) & \(\mathrm m\) & alle vier Koordinaten gleich dimensioniert \\
Eigenzeit \(\tau\) & \(\mathrm s\) & entlang einer zeitartigen Weltlinie \\
Energie \(E\) & \(\mathrm J=\mathrm{kg\,m^2\,s^{-2}}\) & physikalische Observable \\
Wirkung \(S\) & \(\mathrm{J\,s}\) & nicht mit Entropie zu verwechseln \\
Temperatur \(T\) & \(\mathrm K\) & thermodynamische Temperatur \\
Entropie \(S_{\mathrm{th}}\) & \(\mathrm{J\,K^{-1}}\) & \(S_{\mathrm{th}}/k_{\mathrm B}\) dimensionslos \\
Information \(I\) & \(1\) & Bit oder Nat legt die Logarithmusbasis fest \\
Skalarkrümmung \(R\) & \(\mathrm{m^{-2}}\) & bei längendimensionalen Koordinaten \\
Stress-Energie \(T_{\mu\nu}\) & \(\mathrm{J\,m^{-3}}\) & im lokalen orthonormalen Bezugssystem \\
\bottomrule
\end{longtable}

Tensor-Komponenten können bei allgemeinen Koordinatenwechseln andere
Einheitendarstellungen annehmen. Belastbare Aussagen werden deshalb
tensorial oder durch Invarianten formuliert.

\section{Mathematische Metrik und Lorentz-Metrik}
\label{sec:metriken-physik}

Die im Quantisierbarkeitssatz verwendete Metrik
\[
  d:Y\times Y\longrightarrow[0,\infty)
\]
ist positiv definit: \(d(p,q)=0\) gilt genau dann, wenn \(p=q\). Eine
Raumzeitmetrik ist dagegen ein Lorentzsches Tensorfeld \(g_{\mu\nu}\), hier
mit Signatur \((-+++)\). Bei der Koordinatenkonvention \(x^0=ct\) gilt
\[
  \dd s^2=g_{\mu\nu}\,\dd x^\mu\dd x^\nu,
  \qquad [\dd s^2]=\mathrm{m^2},
  \qquad [g_{\mu\nu}]=1.
\]
Für eine zeitartige Kurve ist \(\dd s^2<0\), und ihre Eigenzeit lautet
\[
  \dd\tau=\frac{1}{c}\sqrt{-\dd s^2}.
\]
Auf einer raumartigen Hyperfläche mit induzierter positiv definiter Metrik
\(h_{ij}\) erhält man dagegen \(\dd\ell^2=h_{ij}\dd x^i\dd x^j\)
\citep{wald1984}.

\begin{proposition}[Lorentz-Intervall ist keine Abstandsfunktion]
Das Lorentzsche Linienelement definiert auf der Menge der
Raumzeitereignisse keine Metrik im Sinn eines positiv definiten metrischen
Raums.
\end{proposition}

\begin{proof}
Zwei verschiedene, lichtartig getrennte Ereignisse können \(\dd s^2=0\)
besitzen; für zeitartige Trennungen ist \(\dd s^2<0\). Damit sind positive
Definitheit und die Trennung verschiedener Punkte verletzt.
\end{proof}

Für numerische Konvergenz und endliche \(\eps\)-Netze wird daher eine
zusätzliche positive Darstellungsmetrik benötigt. Die Lorentz-Metrik bestimmt
dagegen Lichtkegel, Eigenzeiten und die kausale Klassifikation von
Trennungen. Beide Objekte heißen ``Metrik'', erfüllen aber verschiedene
mathematische Aufgaben.

\section{Dimensionslose Mandelbrot-Dynamik und Kalibrierung}
\label{sec:nondimensionalisierung}

Die Standarditeration
\[
  z_{n+1}=z_n^2+c
\]
ist dimensionshomogen, wenn \(z_n\) und \(c\) dimensionslose komplexe Zahlen
sind. Wollte man \(z_n\) unmittelbar als dimensionsbehaftete Größe \(Z_n\)
interpretieren, wäre \(Z_{n+1}=Z_n^2+C\) im Allgemeinen unzulässig. Sei
daher \(Q_*>0\) eine Referenzgröße derselben Dimension wie \(Z_n\) und
\(C\). Die dimensionsrichtige Gleichung lautet
\[
  Z_{n+1}=\frac{Z_n^2}{Q_*}+C.
\]
Mit \(z_n=Z_n/Q_*\) und \(c=C/Q_*\) folgt exakt die dimensionslose
Mandelbrot-Iteration. Das Escape-Kriterium \(|z_n|>2\) entspricht in dieser
Skalierung \(|Z_n|>2Q_*\). Die Wahl von \(Q_*\) gehört zum
Korrespondenzmodell und darf nicht verborgen werden.

Sollen mehrere Größen verschiedener Dimensionen in einen komplexen Parameter
eingehen, müssen sie zuerst einzeln normiert werden. Für Observablen \(q_a\)
und Referenzskalen \(q_{a,*}>0\) sei
\[
  \xi_a(X)=\frac{q_a(X)}{q_{a,*}}.
\]
Eine dimensionsrichtige Parameterabbildung hat dann beispielsweise die Form
\[
  \chi(X)=u(\xi_1(X),\ldots,\xi_k(X))
  +\mathrm i\,v(\xi_1(X),\ldots,\xi_k(X))\in\C,
\]
wobei \(u\) und \(v\) dimensionslose reelle Funktionen sind. Geeignete
invariante Argumente können, abhängig vom System, sein:
\[
  \ell_*^2R,\qquad
  \ell_*^4R_{\mu\nu}R^{\mu\nu},\qquad
  \ell_*^4R_{\mu\nu\rho\sigma}R^{\mu\nu\rho\sigma},\qquad
  \frac{E}{E_*},\qquad
  \frac{\tau}{t_*},\qquad
  \frac{S_{\mathrm{th}}}{k_{\mathrm B}}.
\]
Energie und Eigenzeit verlangen zusätzlich die Angabe eines Beobachters oder
einer Weltlinie.

\begin{bemerkung}[Dimensionsrichtigkeit zeichnet keine Abbildung aus]
Bereits
\[
  \chi_{\alpha,\beta}(X)=
  \alpha\ell_*^2R(X)+\mathrm i\,\beta\frac{E(X)}{E_*}
\]
ist für beliebige dimensionslose \(\alpha,\beta\) dimensionsrichtig. Daraus
folgt weder, dass diese Abbildung physikalisch ausgezeichnet ist, noch dass
ihre Mandelbrot-Klassifikation eine Naturklasse beschreibt. Dafür sind
Symmetrien, Dynamik, Observablen und Tests erforderlich.
\end{bemerkung}

Auch der Iterationsindex \(n\) ist dimensionslos. Eine Gleichsetzung mit
physikalischer Zeit verlangt eine weitere Definition,
\[
  \tau_N=\sum_{n=0}^{N-1}\Delta\tau_n,
  \qquad [\Delta\tau_n]=\mathrm s.
\]
Insbesondere folgt aus einer diskreten Iterationsfolge kein fundamentaler
Zeittakt. Die Festlegung \(\Delta\tau_n=t_{\mathrm P}\) wäre eine zusätzliche
Hypothese, keine Konsequenz des Algorithmus.

\subsection{Positive Metrik auf dem Beobachtungsraum}

Für Mess- und Konvergenzfragen kann ein dimensionsloses Darstellungsmaß
definiert werden:
\[
  d_*(X,Y)^2=
  \sum_{a=1}^{k}w_a
  \left(\frac{q_a(X)-q_a(Y)}{q_{a,*}}\right)^2,
  \qquad w_a>0.
\]
Es ist eine Metrik, wenn die ausgewählten Observablen die betrachteten
Zustände trennen; andernfalls ist es nur eine Pseudometrik. Aus
\(d_*(X,Y)\le\eps\) folgt für jede Komponente
\[
  |q_a(X)-q_a(Y)|\le
  \frac{\eps}{\sqrt{w_a}}q_{a,*}.
\]
Damit wird die abstrakte Genauigkeit \(\eps\) in konkrete Auflösungen in
Metern, Sekunden oder Joule übersetzt. Die Wahl von \(q_{a,*}\) und \(w_a\)
ist eine dokumentierte Mess- und Modellkonvention, kein Beweis fundamentaler
Körnung.

Für mehr als zwei reelle Freiheitsgrade ist ein einzelnes \(c\in\C\) im
Allgemeinen nicht injektiv. Eine höherdimensionale physikalische Theorie
benötigt deshalb eine Familie von Parametern, etwa einen Zustand in \(\C^d\),
oder eine offengelegte Merkmalsprojektion samt quantifiziertem
Informationsverlust. Dies konkretisiert die Bedingung der
Klassentrennung aus \cref{satz:abbildung,satz:gate-faktorisierung}.

\section{Physikalische Typisierung der Wirkraumfaktoren und des Gates}
\label{sec:rtcia-typen}

Die Kurzform
\[
  \mathcal W=\mathcal A
  (\mathcal R\times\mathcal T\times\mathcal C_{\mathrm{caus}}\times\mathcal I)
\]
ist dimensionsrichtig zu lesen, wenn \(\times\) ein kartesisches Produkt
typisierter Strukturen und keine Multiplikation von Zahlenwerten bezeichnet.
Formal ist
\[
  \mathcal A:\mathfrak R\times\mathfrak T\times\mathfrak C\times
  \mathfrak I\longrightarrow\mathfrak S
\]
eine Bewertungs- oder Freigabeabbildung in einen Statusraum \(\mathfrak S\).

\begin{longtable}{>{\raggedright\arraybackslash}p{0.15\textwidth}
                  >{\raggedright\arraybackslash}p{0.31\textwidth}
                  >{\raggedright\arraybackslash}p{0.43\textwidth}}
\toprule
\textbf{Symbol} & \textbf{Mathematischer Typ} & \textbf{Physikalische Einheit oder Bedeutung} \\
\midrule
\endfirsthead
\toprule
\textbf{Symbol} & \textbf{Mathematischer Typ} & \textbf{Physikalische Einheit oder Bedeutung} \\
\midrule
\endhead
\(\mathcal R\) & Relation, Metrik- oder Geometriefeld & keine einzelne SI-Einheit; Längen in \(\mathrm m\), Krümmungen in \(\mathrm{m^{-2}}\) \\
\(\mathcal T\) & Ordnungsrelation oder Folgenindex & \(n\) dimensionslos; Eigenzeit und Zeitabstand in \(\mathrm s\) \\
\(\mathcal C_{\mathrm{caus}}\) & Prädikat, Übergangskern oder Operator & Boolescher Status und Wahrscheinlichkeit dimensionslos; Zeitgenerator gegebenenfalls \(\mathrm{s^{-1}}\) \\
\(\mathcal I\) & Informationsmaß oder Evidenzstruktur & Shannon-Information dimensionslos in Bit oder Nat; thermodynamische Entropie gegebenenfalls \(S_{\mathrm{th}}=k_{\mathrm B}I\) in \(\mathrm{J\,K^{-1}}\) \citep{shannon1948a} \\
\(\mathcal A\) & Gate- oder Prüfabbildung & keine SI-Einheit; liefert Status oder typisierte Wirkung \\
\(\mathcal W\) & akzeptierter Zustand oder Wirkung & Einheit hängt vom Zieltyp ab; nicht automatisch die SI-Einheit Watt \\
\bottomrule
\end{longtable}

Auch in einer Eigenwertgleichung
\[
  \widehat{\mathcal C}\lvert\psi\rangle
  =\lambda\lvert\psi\rangle
\]
wird die Einheit von \(\lambda\) durch den Operatortyp festgelegt. Für einen
Projektor ist \(\lambda\in\{0,1\}\) dimensionslos, für einen Hamiltonoperator
ist \(\lambda\) eine Energie in Joule, und für einen Zeitgenerator kann
\(\lambda\) die Einheit \(\mathrm{s^{-1}}\) besitzen. Eine universale
``Wirkstärke'' ist deshalb erst nach Angabe des Operators bestimmt.

\section{Geometrie, Materie und Quantenphase}
\label{sec:physikalische-kopplung}

\subsection{Einstein-Gleichung und kausale Struktur}

Mit der festgelegten Koordinatenkonvention lauten die Einstein-Gleichungen
\citep{einstein1916,wald1984}
\[
  G_{\mu\nu}+\Lambda g_{\mu\nu}
  =\frac{8\pi G}{c^4}T_{\mu\nu}.
\]
Im lokalen orthonormalen Bezugssystem gilt
\[
  [G_{\mu\nu}]=[\Lambda]=\mathrm{m^{-2}},
  \qquad [T_{\mu\nu}]=\mathrm{J\,m^{-3}},
\]
und
\[
  \left[\frac{G}{c^4}T_{\mu\nu}\right]=\mathrm{m^{-2}}.
\]
Die Gleichung ist also dimensionshomogen. Ihre physikalische Aussage geht
darüber hinaus: Materie- und Feldzustände tragen Stress-Energie; die
gekoppelte Geometrie bestimmt ihrerseits Lichtkegel, Eigenzeiten und die
kovariante Ausbreitung der Materiefelder. Die Einstein-Materie-Theorie ist
daher ein gekoppeltes Anfangswertproblem, nicht bloß eine einseitige Kette
``Materie verursacht Geometrie''.

Die geometrische Bianchi-Identität
\[
  \nabla^\mu G_{\mu\nu}=0
\]
führt zusammen mit den Einstein-Gleichungen bei konstantem \(\Lambda\) zu
\[
  \nabla^\mu T_{\mu\nu}=0.
\]
Eine physikalische Anschlussregel, die diese lokale kovariante
Energie-Impuls-Erhaltung verletzt, wäre nicht ohne zusätzliche neue Dynamik
mit der Allgemeinen Relativitätstheorie vereinbar.

\subsection{Wirkungsprinzip und Quantenskala}

Für \(x^0=ct\) kann die Einstein--Hilbert-Wirkung geschrieben werden als
\[
  S_{\mathrm{EH}}[g,\Psi]
  =\frac{c^3}{16\pi G}
  \int(R-2\Lambda)\sqrt{-g}\,\dd^4x
  +S_{\mathrm m}[g,\Psi].
\]
Da \([\dd^4x]=\mathrm{m^4}\), \([R]=\mathrm{m^{-2}}\) und
\([c^3/G]=\mathrm{kg\,s^{-1}}\), besitzt \(S_{\mathrm{EH}}\) die Einheit
\[
  [S_{\mathrm{EH}}]=\mathrm{kg\,m^2\,s^{-1}}=\mathrm{J\,s}.
\]
Die Variation nach \(g_{\mu\nu}\) und den Materiefeldern \(\Psi\) ergibt
die gekoppelten Feldgleichungen.

In einer quantenmechanischen Beschreibung tritt die Wirkung in der
dimensionslosen Phase
\[
  \exp\!\left(\frac{\mathrm iS}{\hbar}\right)
\]
auf \citep{feynman1948}. Hier vermittelt \(\hbar\) zwischen Wirkung und
Quantenphase. Entsprechend sind
\[
  E=\hbar\omega,
  \qquad p=\hbar k,
\]
mit \([\omega]=\mathrm{s^{-1}}\) und \([k]=\mathrm{m^{-1}}\)
dimensionshomogen.

Damit ist ein präziser Standardanschluss gegeben: Die Lorentz-Metrik liefert
geometrische und kausale Struktur, \(T_{\mu\nu}\) repräsentiert
Energie-Impuls, das Wirkungsprinzip erzeugt die gekoppelte Dynamik, und
\(S/\hbar\) bildet die dimensionslose Größe der Quantenphase. Daraus folgt
noch keine bestimmte Theorie der Quantengravitation.

\section{Planck-Einheiten: natürliche Skalen, keine bewiesenen Gitterabstände}
\label{sec:planck-einheiten}

Aus \(\hbar\), \(G\) und \(c\) entstehen die dimensionsrichtigen
Planck-Einheiten \citep{codata2022}:
\[
\begin{aligned}
  \ell_{\mathrm P}&=\sqrt{\frac{\hbar G}{c^3}},
  & [\ell_{\mathrm P}]&=\mathrm m,\\
  t_{\mathrm P}&=\sqrt{\frac{\hbar G}{c^5}}
  =\frac{\ell_{\mathrm P}}{c},
  & [t_{\mathrm P}]&=\mathrm s,\\
  m_{\mathrm P}&=\sqrt{\frac{\hbar c}{G}},
  & [m_{\mathrm P}]&=\mathrm{kg},\\
  E_{\mathrm P}&=m_{\mathrm P}c^2
  =\sqrt{\frac{\hbar c^5}{G}},
  & [E_{\mathrm P}]&=\mathrm J.
\end{aligned}
\]
In ihnen treten Quantenwirkung und Gravitation sowie die relativistische
Kausalgeschwindigkeit gemeinsam auf. Sie liefern natürliche Normierungen
\[
  \widetilde\ell=\frac{\ell}{\ell_{\mathrm P}},
  \qquad
  \widetilde\tau=\frac{\tau}{t_{\mathrm P}},
  \qquad
  \widetilde E=\frac{E}{E_{\mathrm P}}.
\]
Nicht daraus folgt dagegen
\[
  \ell=n\ell_{\mathrm P},
  \qquad \tau=nt_{\mathrm P}
  \qquad(n\in\mathbb Z)
\]
für alle Längen oder Zeiten. Eine solche Spektral- oder Gitteraussage
benötigt einen definierten Operator, ein Zustandsspektrum und messbare
Konsequenzen. Die Planck-Einheiten sind charakteristische Skalen, keine
bereits nachgewiesenen Raumzeitatome \citep{hossenfelder2013}.

\section{Schwarzloch-Thermodynamik als physikalische Brücke}
\label{sec:horizont-information}

Für ein ungeladenes, nichtrotierendes Schwarzschild-Loch gilt
\[
  r_{\mathrm s}=\frac{2GM}{c^2},
  \qquad A_{\mathrm H}=4\pi r_{\mathrm s}^2.
\]
Die Bekenstein--Hawking-Entropie lautet \citep{bekenstein1973,hawking1975}
\[
  S_{\mathrm{BH}}
  =\frac{k_{\mathrm B}c^3A_{\mathrm H}}{4G\hbar}
  =\frac{k_{\mathrm B}A_{\mathrm H}}{4\ell_{\mathrm P}^2},
  \qquad [S_{\mathrm{BH}}]=\mathrm{J\,K^{-1}}.
\]
Der dimensionslose Entropiezähler in Nat und in Bit ist
\[
  I_{\mathrm{BH}}^{(\mathrm{nat})}
  =\frac{S_{\mathrm{BH}}}{k_{\mathrm B}}
  =\frac{A_{\mathrm H}}{4\ell_{\mathrm P}^2},
  \qquad
  I_{\mathrm{BH}}^{(\mathrm{bit})}
  =\frac{S_{\mathrm{BH}}}{k_{\mathrm B}\ln2}.
\]
Damit sind zwei Oberflächendichten zu unterscheiden:
\[
  \frac{\dd S_{\mathrm{BH}}}{\dd A_{\mathrm H}}
  =\frac{k_{\mathrm B}}{4\ell_{\mathrm P}^2}
  \quad\left[\frac{\mathrm J}{\mathrm{K\,m^2}}\right],
\]
und
\[
  \frac{\dd I_{\mathrm{BH}}^{(\mathrm{nat})}}{\dd A_{\mathrm H}}
  =\frac{1}{4\ell_{\mathrm P}^2}
  \quad[\mathrm{m^{-2}}].
\]
Ein Ausdruck \(I_{\mathrm{horizon}}=\Delta S/\Delta A\) bezeichnet daher,
wenn \(S\) thermodynamische Entropie ist, zunächst \emph{Entropie pro
Fläche}. Erst die Division durch \(k_{\mathrm B}\) ergibt eine
dimensionslose Zustandszählung pro Fläche.

Die Hawking-Temperatur ist
\[
  T_{\mathrm H}=\frac{\hbar c^3}{8\pi GMk_{\mathrm B}}.
\]
Für die Schwarzschild-Familie gilt konsistent
\[
  \dd(Mc^2)=T_{\mathrm H}\,\dd S_{\mathrm{BH}}.
\]
Hier treffen Geometrie \((A_{\mathrm H})\), Gravitation \((G)\),
Quantentheorie \((\hbar)\) und Thermodynamik
\((k_{\mathrm B},T_{\mathrm H})\) in einer quantitativ bestimmten Relation
zusammen. Dies ist eine zentrale Brücke, aber noch keine vollständige
mikroskopische Quantengravitation.

Auch das Landauer-Prinzip ist eng zu begrenzen: Für die logisch irreversible
Löschung eines Bits in einem thermischen Umfeld bei Temperatur \(T\) gilt im
idealen quasistatischen Grenzfall
\[
  Q_{\min}=k_{\mathrm B}T\ln2
\]
\citep{landauer1961}. Die Formel behauptet weder, dass jedes Bit diese
Energie besitzt, noch dass abstrakte Information unabhängig von einem
physischen Träger gravitiert. In der etablierten Allgemeinen
Relativitätstheorie koppelt der Stress-Energie-Tensor des Trägers an die
Geometrie, nicht ein semantischer Informationsgehalt als selbständige
Quellenart.

\section{Ein physikalisch typisiertes Anschluss- und Gatemodell}
\label{sec:phys-gate}

Die Anschlussidee kann als prüfbare numerische Architektur formuliert werden,
ohne sie bereits mit einer fundamentalen Naturdynamik gleichzusetzen. Sei
\(\Sigma_n\) eine Folge raumartiger Hyperflächen und
\[
  X_n=\bigl(h_{ij}^{(n)},K_{ij}^{(n)},\Psi_n,\Pi_n;D_n\bigr)
\]
der Zustandsdatensatz aus induzierter räumlicher Metrik \(h_{ij}\),
extrinsischer Krümmung \(K_{ij}\), Materiefeldern \(\Psi\), ihren kanonischen
Daten \(\Pi\) und Messdaten \(D_n\). Ein Evolutionsoperator erzeugt zunächst
nur einen Kandidaten:
\[
  \widehat X_{n+1}=G_{\Delta\tau_n}(X_n),
  \qquad [\Delta\tau_n]=\mathrm s.
\]
Ein physikalisches Gate prüft anschließend
\[
  \mathcal A_{\mathrm{phys}}:
  (X_n,\widehat X_{n+1},D_n)
  \longrightarrow\{\PASS,\CONTINUE,\BLOCK\}.
\]

Seien \(C_a(\widehat X_{n+1})\) die Residuen der relevanten
Zwangs-, Erhaltungs- und Evolutionsgleichungen und \(C_{a,*}>0\) vorab
festgelegte Referenzgrößen derselben jeweiligen Dimension. Dann sind
\[
  r_a:=\frac{\lVert C_a(\widehat X_{n+1})\rVert}{C_{a,*}}
\]
dimensionslos. Zusammen mit normierten Messunsicherheiten \(u_b\), einer
Toleranz \(\eta\) und einem separaten Test der durch \(g_{\mu\nu}\)
bestimmten kausalen Abhängigkeit kann ein schematisches Gate lauten:
\[
\mathcal A_{\mathrm{phys}}=\PASS
\Longleftrightarrow
\max\{r_a,u_b\}\le\eta
\quad\text{und alle notwendigen Kausalbedingungen sind erfüllt}.
\]
Die Größen \(C_{a,*}\), die Normen, \(\eta\) und die geprüften Bedingungen
müssen vor der Auswertung dokumentiert werden; andernfalls wäre das Gate
nicht reproduzierbar.

Der freigegebene Schritt lautet
\[
  X_{n+1}=\operatorname{Commit}
  \bigl(\widehat X_{n+1}\mid
  \mathcal A_{\mathrm{phys}}=\PASS\bigr).
\]
Dies ist eine wissenschaftlich sinnvolle Wirkungshaltepunkt-Architektur: Ein
berechneter Zustand gilt nicht allein aufgrund seiner Berechenbarkeit als
physikalisch zulässig.

Ein physikalisch relevanter Fixpunkt \(X^*\) muss mindestens erfüllen
\[
  G_{\Delta\tau}(X^*)=X^*,
  \qquad
  \mathcal A_{\mathrm{phys}}(X^*,X^*,D)=\PASS.
\]
Führt ein blockierter Kandidat lediglich dazu, dass der alte Zustand
unverändert gespeichert bleibt, entsteht ebenfalls ein administrativer
Fixpunkt. Dieser ist kein Nachweis einer stationären Lösung:
\[
  \boxed{
  \text{Unverändertheit durch BLOCK}
  \neq
  \text{physikalisch freigegebener Fixpunkt}.}
\]

Werden Raum und Zeit numerisch mit Gitterweiten \(h\) und \(\Delta\tau\)
diskretisiert, verlangt die Validierung eine Konvergenzaussage der Form
\[
  \lim_{h,\Delta\tau\to0}
  \bigl\|X^{(h,\Delta\tau)}-X\bigr\|_*=0
\]
sowie Stabilitäts- und Konsistenztests. Solange sich das Raster im
Grenzübergang verfeinern lässt, ist es eine Darstellungsmethode. Ein
fundamentales Raumzeitquant wäre erst durch ein Spektrum physikalischer
Observablen und messbare Konsequenzen begründet.

\section{Explizite physikalische Korrespondenz zur Mandelbrot-Dynamik}
\label{sec:phys-semikonjugation}

Für eine dynamische physikalische Interpretation genügt eine bloße Zuordnung
\(X\mapsto c\) nicht. Sei \(\mathfrak X\) ein physikalischer Zustandsraum,
\(G_{\Delta\tau}:\mathfrak X\to\mathfrak X\) die Zeitentwicklung und
\[
  \Phi:\mathfrak X\longrightarrow Y_U
\]
eine aus dimensionslosen, möglichst koordinateninvarianten Observablen
gebildete Abbildung. Eine exakte dynamische Korrespondenz verlangt mindestens
\[
  \boxed{\Phi\circ G_{\Delta\tau}=F\circ\Phi.}
\]
Diese Richtung kodiert physikalische Zustände im Modell. Sie ist von der in
Teil II betrachteten Realisierungsrichtung
\(G\circ\phi=\phi\circ F\) zu unterscheiden. Existieren beide Abbildungen
als zueinander inverse Bijektionen, liegt eine Konjugation vor; eine bloße
Semikonjugation darf Informationen verlieren.

Für eine angenäherte Korrespondenz wird ein vorab normiertes Residuum
benötigt:
\[
  \delta_\Phi(X)=d_{Y_U}
  \bigl(\Phi(G_{\Delta\tau}X),F(\Phi(X))\bigr),
  \qquad \delta_\Phi(X)\le\eta_\Phi.
\]
Die positive Metrik \(d_{Y_U}\), die Toleranz \(\eta_\Phi\) und der
Gültigkeitsbereich sind vor dem Test festzulegen.

Soll die Mandelbrot-Klassifikation eine physikalische
Fortsetzbarkeitsklasse darstellen, ist zusätzlich Gate-Verträglichkeit nötig:
\[
  \mathcal A_{\mathrm{phys}}(X)
  =\mathcal A_{\M}(\Phi(X))
\]
für die beanspruchte Klasse von Zuständen. Diese Bedingung ist wesentlich
stärker als eine visuelle Ähnlichkeit von Fraktalen und physikalischen
Strukturen. Die Hypothese zerfällt daher in drei unabhängige Prüfaufgaben:

\begin{enumerate}
  \item \textbf{Dimension und Invarianz:} \(\Phi\) besteht aus
  wohldefinierten dimensionslosen Observablen und ist nicht bloß ein Effekt
  willkürlicher Koordinatenwahl.
  \item \textbf{Dynamik:} Semikonjugation oder quantitativ begrenztes
  Residuum gelten für Lösungen der physikalischen Gleichungen.
  \item \textbf{Klassifikation:} PASS-, CONTINUE- und BLOCK-Klassen stimmen
  mit einer unabhängig messbaren physikalischen Eigenschaft überein.
\end{enumerate}

Ohne diese Nachweise bleibt die Mandelbrot-Struktur ein mathematisch
anschlussfähiges Modell und eine Heuristik. Mit ihnen würde sie zu einer
konkret prüfbaren physikalischen Korrespondenzhypothese.

\section{Was beruht auf was? Logische, metrische und kausale Pfeile}
\label{sec:kausalkette}

Die Pfeile des Gesamtmodells dürfen nicht sämtlich als physikalische Ursachen
gelesen werden. Mindestens fünf Relationstypen sind zu unterscheiden:

\begin{longtable}{>{\raggedright\arraybackslash}p{0.22\textwidth}
                  >{\raggedright\arraybackslash}p{0.29\textwidth}
                  >{\raggedright\arraybackslash}p{0.38\textwidth}}
\toprule
\textbf{Relation} & \textbf{Beispiel} & \textbf{Bedeutung} \\
\midrule
\endfirsthead
\toprule
\textbf{Relation} & \textbf{Beispiel} & \textbf{Bedeutung} \\
\midrule
\endhead
transzendentale Voraussetzung & Unterschied \(\to\) Bestimmbarkeit & ohne unterscheidbare Zustände keine bestimmte Aussage; kein zeitlicher Naturprozess behauptet \\
metrologische Repräsentation & Größe \(q\), Einheit \(u_q\), Zahlenwert & Kalibrierung erzeugt eine Darstellung, nicht das gemessene Phänomen \\
konstitutive Geometrie & \(g_{\mu\nu}\to\) Eigenzeiten und Lichtkegel & die Metrik definiert im Modell räumlich-zeitliche und kausale Relationen \\
dynamische Kopplung & \((g_{\mu\nu},T_{\mu\nu})\to\) Entwicklung & Einstein- und Materiegleichungen bestimmen aus Anfangsdaten zulässige Entwicklungen \\
epistemische oder technische Freigabe & Kandidat und Evidenz \(\to\) Status & das Gate verursacht die Freigabe einer berechneten Wirkung, nicht den zugrunde liegenden Naturprozess \\
\bottomrule
\end{longtable}

Die physikalisch belastbare Kette kann deshalb so geschrieben werden:
\[
\boxed{
\begin{gathered}
\text{Wirkung }S[g,\Psi]+\text{ Anfangs- und Randdaten}
\\
\Downarrow\ \text{Variation und Feldgleichungen}
\\
(g_{\mu\nu},\Psi,T_{\mu\nu})
\\
\Downarrow\ \text{kausale Dynamik}
\\
\text{physikalische Observablen }q_a
\\
\Downarrow\ \text{Messung und Normierung}
\\
\xi_a=q_a/q_{a,*}
\\
\Downarrow\ \text{definierte Korrespondenz }\Phi
\\
\text{dimensionsloser Modellzustand}
\\
\Downarrow\ \text{Iteration und Gate}
\\
\text{Modellklassifikation und freigegebene Folgewirkung}.
\end{gathered}}
\]

Die SI-Einheiten stehen auf der Seite der Messung und Kalibrierung. Sie
verursachen weder Raumzeit noch Energie. Die Metrik legt innerhalb der Theorie
fest, welche Ereignisse kausal verbunden sein können. Die gekoppelte Dynamik
und die Anfangs- und Randdaten legen fest, welche Entwicklung modelliert
wird. Ein QIK--VRT-Gate kann anschließend steuern, ob ein berechnetes Ergebnis
als technische Wirkung freigegeben wird.

Für Information gilt entsprechend
\[
  \text{abstrakter Informationsgehalt}
  \neq\text{physischer Träger}
  \neq T_{\mu\nu}\text{ des Trägers}.
\]
Dass ein physischer Speicher oder ein Signal gravitiert, folgt aus seinem
Energie-Impuls. Eine zusätzliche direkte Kopplung semantischer Information an
die Raumzeit ist in der etablierten Einstein-Theorie nicht enthalten und
müsste als neue Dynamik formuliert und getestet werden. Ebenso ist die
spätere Neuklassifikation früherer Iterationswerte keine physikalische
Rückwärtskausalität. Sie verändert den epistemischen Status einer Spur, nicht
deren frühere Zustandswerte.

\clearpage
\part{Retrokausalität: Begriffe, Modelle, Tests und Konsequenzen}

\section{Retrokausalität: Begriffe, Modelle und empirische Grenzen}
\label{sec:retrokausalitaet}

Der Ausdruck \emph{Retrokausalität} wird in Physik, Philosophie,
Quantenfundamenten und populärwissenschaftlicher Darstellung in mehreren
verschiedenen Bedeutungen verwendet. Ohne eine explizite Typisierung kann
derselbe Ausdruck eine harmlose nachträgliche Neubewertung, eine
zeitumkehrsymmetrische Gleichung, eine Zwei-Randwertformulierung, eine
ontologische Abhängigkeit von späteren Bedingungen oder sogar einen
kontrollierbaren Nachrichtenkanal in die Vergangenheit bezeichnen. Diese
Bedeutungen sind weder logisch noch empirisch gleichwertig.

Der vorliegende Abschnitt trennt deshalb drei Leitfragen:
\begin{enumerate}
  \item \emph{Was wird mathematisch oder begrifflich behauptet?}
  \item \emph{Welche physikalische Abhängigkeit würde daraus folgen?}
  \item \emph{Welche Beobachtung könnte diese Abhängigkeit von einer bloßen
  Korrelation, Nachselektion oder gewöhnlichen Vorwärtskausalität
  unterscheiden?}
\end{enumerate}
Die stärkste wissenschaftlich belastbare Kurzform lautet
\[
\boxed{
\begin{gathered}
\text{Rückschluss}\neq\text{Rückwirkung},\\
\text{Zeitumkehrsymmetrie}\neq\text{Retrokausalität},\\
\text{retrokausale Modellierung}\neq
\text{rückwärts gerichtete Signalisierung}.
\end{gathered}}
\]

\subsection{Minimaldefinitionen: Zustand, Evidenz, Intervention und Signal}

Seien \(t_1<t_2\) zwei physikalisch geordnete Zeitpunkte. Wir unterscheiden:
\begin{itemize}
  \item \(X_{t_1}\): einen früheren ontischen oder modellierten Zustand;
  \item \(R_{t_1}\): einen zum Zeitpunkt \(t_1\) erzeugten und gegen spätere
  Veränderung gesicherten Mess- oder Datenrecord;
  \item \(A_{t_2}\): eine später gewählte experimentelle Einstellung;
  \item \(Y_{t_2}\): ein späteres Messergebnis;
  \item \(\Lambda\): offengelegte Vorbedingungen und mögliche gemeinsame
  Ursachen.
\end{itemize}

Eine statistische Abhängigkeit
\[
  \mathbb P(X_{t_1}\mid A_{t_2}=a)
  \neq
  \mathbb P(X_{t_1}\mid A_{t_2}=a')
\]
beweist für sich genommen keine Kausalität. Sie kann durch gemeinsame
Ursachen, Auswahlverzerrung, Datenleckage, Nachselektion oder eine
unvollständige Modellierung entstehen. Der interventionistische Unterschied
zwischen Beobachtung und Eingriff wird durch den
\(\operatorname{do}\)-Operator markiert \citep{Pearl2009}. Eine
\emph{retrokausale Zustandsabhängigkeit} läge in einem bestimmten Modell vor,
wenn für \(t_1<t_2\) und sonst festgehaltene Randbedingungen
\begin{equation}
\label{eq:retro-ontisch}
  \mathbb P\!\left(X_{t_1}\mid
  \operatorname{do}(A_{t_2}=a),\Lambda\right)
  \neq
  \mathbb P\!\left(X_{t_1}\mid
  \operatorname{do}(A_{t_2}=a'),\Lambda\right)
\end{equation}
gälte und die Differenz nicht über einen gewöhnlichen vorwärts gerichteten
Pfad oder eine gemeinsame Ursache vermittelt wäre.

Davon nochmals verschieden ist eine \emph{operative rückwärts gerichtete
Signalisierung}. Sie verlangt, dass ein frei kontrollierter späterer Eingriff
die Verteilung eines früher zugänglichen Records verändert:
\begin{equation}
\label{eq:retro-signal}
  \mathbb P\!\left(R_{t_1}=r\mid
  \operatorname{do}(A_{t_2}=a)\right)
  \neq
  \mathbb P\!\left(R_{t_1}=r\mid
  \operatorname{do}(A_{t_2}=a')\right).
\end{equation}
Erst wenn \(a\) als Nachricht codiert und aus \(R_{t_1}\) mit einer über
Zufall hinausgehenden Zuverlässigkeit vor \(t_2\) decodiert werden könnte,
bestünde ein nutzbarer Rückwärtskanal. Eine ontologische Interpretation kann
also retrokausal sein, ohne \cref{eq:retro-signal} zu erfüllen.

\begin{definition}[Semantische Rückbestimmung]
Eine semantische Rückbestimmung liegt vor, wenn spätere Evidenz die zulässige
Klassifikation oder Bedeutung einer früher erzeugten Spur verändert, ohne
deren damaligen physischen Zustand oder ihren gesicherten Record zu ändern.
\end{definition}

\begin{definition}[Physikalische Retrokausalität]
Physikalische Retrokausalität bezeichnet hier eine modellierte oder reale
kausale Abhängigkeit eines früheren physischen Zustands von einer späteren
Randbedingung, Einstellung oder einem späteren Ereignis, die nicht vollständig
auf gewöhnliche Vorwärtskausalität, gemeinsame Ursachen oder
Konditionierung zurückgeführt werden kann.
\end{definition}

\begin{definition}[Rückwärtssignalisierung]
Rückwärtssignalisierung ist die kontrollierbare Übertragung einer frei
wählbaren Nachricht zu einem früheren, dort bereits auslesbaren Record. Sie
ist stärker als eine retrokausale Ontologie und stärker als eine
zeit-symmetrische mathematische Beschreibung.
\end{definition}

Diese Definitionen verhindern zwei symmetrische Fehlschlüsse. Aus dem Fehlen
eines Rückwärtskanals folgt nicht, dass jede zeit-symmetrische oder
retrokausale Interpretation widersprüchlich wäre. Aus der mathematischen
Konsistenz einer solchen Interpretation folgt umgekehrt kein experimenteller
Nachweis einer Wirkung in die Vergangenheit.

\subsection{Eine Taxonomie zeitlicher Rückbezüge}

Die folgende Tabelle ordnet die gebräuchlichsten Bedeutungen nach dem Objekt,
das sich tatsächlich ändert.

\begin{longtable}{>{\raggedright\arraybackslash}p{0.22\textwidth}
                  >{\raggedright\arraybackslash}p{0.29\textwidth}
                  >{\raggedright\arraybackslash}p{0.39\textwidth}}
\toprule
\textbf{Begriff} & \textbf{Was ändert sich?} & \textbf{Was folgt daraus nicht?} \\
\midrule
\endfirsthead
\toprule
\textbf{Begriff} & \textbf{Was ändert sich?} & \textbf{Was folgt daraus nicht?} \\
\midrule
\endhead
semantische Rückbestimmung & Klassifikation einer unveränderten früheren Spur & keine Änderung des früheren Zustands oder Records \\
Retrodiktion oder Smoothing & Wahrscheinlichkeitsurteil über einen früheren verborgenen Zustand & keine physische Rückwirkung der späteren Daten \\
Nachselektion & Zusammensetzung des später betrachteten Teilensembles & keine Änderung der unbedingten früheren Stichprobe \\
Zeitumkehrinvarianz & Form der Gleichung unter einer Zeitumkehrtransformation & keine Aussage, welche Randbedingung physisch realisiert ist \\
Zwei-Randwertproblem & zulässige Geschichte wird durch Anfangs- und Enddaten eingeschränkt & kein frei kontrollierbarer Rückwärtskanal \\
avancierte Lösung & mathematische Lösung mit Unterstützung aus der zeitlichen Zukunft & keine automatische physikalische Auswahl dieser Lösung \\
ontische Retrokausalität & früherer Modellzustand hängt von späterer Bedingung ab & keine automatische Beobachtbarkeit oder Signalisierung \\
Rückwärts\-signalisierung & früherer zugänglicher Record trägt eine später frei gewählte Nachricht & keine etablierte Konsequenz der Standardquantenmechanik \\
geschlossene zeitartige Kurve & Weltlinie kehrt in ihre eigene kausale Vergangenheit zurück & keine Gleichsetzung mit bloßer Vor-/Nachselektion \\
indefinite Kausalordnung & Prozesse besitzen keine feste globale Reihenfolge bestimmter Operationen & nicht notwendig eine Wirkung in die eigene Vergangenheit \\
\bottomrule
\end{longtable}

Besonders häufig wird Konditionierung mit Kausalität verwechselt. Für einen
früheren Zustand \(X\), frühere Daten \(D_1\) und spätere Daten \(D_2\) gilt
nach Bayes
\begin{equation}
\label{eq:smoothing}
 \mathbb P(X\mid D_1,D_2)
 =\frac{\mathbb P(D_2\mid X,D_1)\,
              \mathbb P(X\mid D_1)}
             {\mathbb P(D_2\mid D_1)}.
\end{equation}
Der Übergang von \(\mathbb P(X\mid D_1)\) zu
\(\mathbb P(X\mid D_1,D_2)\) verändert ein Informationsurteil. Er schreibt
keine neue Bewegungsgleichung für den bereits vergangenen Zustand. In
Signalverarbeitung, Navigation und Zustandsrekonstruktion ist dieses
\emph{Smoothing} alltäglich und nützlich; gerade deshalb darf es nicht als
experimenteller Nachweis physikalischer Retrokausalität ausgegeben werden.

\subsection{Relativistische Kausalstruktur und Zeitordnung}

In einer zeit-orientierten Lorentz-Mannigfaltigkeit
\((\mathcal X,g_{\mu\nu})\) bezeichnet \(J^+(p)\) die Menge der durch
zukünftig gerichtete zeit- oder lichtartige Kurven von \(p\) erreichbaren
Ereignisse; \(J^-(p)\) bezeichnet die entsprechende kausale Vergangenheit.
Die Kausalrelation ist damit nicht eine beliebige Pfeilrichtung in einem
Diagramm, sondern wird durch die Lichtkegelstruktur der Metrik bestimmt
\citep{wald1984}.

Für zwei Ereignisse \(p,q\) sind mindestens drei Fälle zu unterscheiden:
\begin{align*}
 q\in I^+(p) &\quad\text{zeitartig erreichbar},\\
 q\in J^+(p)\setminus I^+(p) &\quad\text{lichtartig erreichbar},\\
 q\notin J^+(p)\cup J^-(p) &\quad\text{raumartig getrennt}.
\end{align*}
Die Reihenfolge raumartig getrennter Ereignisse kann von der Wahl des
Inertialsystems abhängen. Daraus folgt keine physische Rückwirkung: Eine
invariante Signalaussage verlangt einen kausalen Übertragungspfad. Umgekehrt
kann eine hypothetische kontrollierbare Überlichtgeschwindigkeit zusammen mit
Lorentz-Symmetrie in geeigneten Bezugssystemen zu kausalen Schleifen führen.
Deshalb sind Überlichtsignalisierung, Retrokausalität und geschlossene
zeitartige Kurven eng verwandt, aber nicht definitorisch identisch.

Auf einer global hyperbolischen Raumzeit kann eine Cauchyfläche so gewählt
werden, dass geeignete hyperbolische Feldgleichungen aus Anfangsdaten eine
wohlgestellte Entwicklung besitzen. Dies ist die Standardstruktur eines
Vorwärtsproblems. Eine mathematische Theorie kann stattdessen Randdaten an
mehreren Zeiten verwenden. Dann wird eine ganze Geschichte als Lösung eines
Randwertproblems bestimmt. Das Wort ``bestimmt'' bezeichnet zunächst eine
globale mathematische Einschränkung; es legt noch nicht fest, welcher Rand als
wirkende Ursache und welcher als Wirkung interpretiert wird.

\begin{bemerkung}[Koordinaten- und Kausalrichtung]
Ein negativer Koordinaten\-zeitschritt, eine rückwärts parametrisierte Kurve oder
das formale Invertieren eines Evolutionsoperators beweist keine
Retrokausalität. Physikalisch relevant sind invariant definierte
Kausalbeziehungen, Randbedingungen und die Möglichkeit operationaler
Interventionen.
\end{bemerkung}

\subsection{Zeitumkehrsymmetrie, Bewegungsgesetz und thermodynamischer Pfeil}

Sei eine Dynamik durch einen Fluss \(U_t\) auf einem Zustandsraum beschrieben.
Eine Zeitumkehroperation \(\Theta\) erfüllt in einem zeitumkehrinvarianten
Modell beispielsweise
\begin{equation}
\label{eq:time-reversal}
  \Theta U_t\Theta^{-1}=U_{-t}.
\end{equation}
Diese Gleichung besagt: Zu einer zulässigen Geschichte gehört unter
geeigneter Transformation der zeitungeraden Größen eine zeitumgekehrte
zulässige Geschichte. Sie besagt nicht, dass in einer gegebenen realisierten
Geschichte spätere freie Entscheidungen frühere Zustände erzeugen.
Zeitumkehrinvarianz ist zudem keine universelle Eigenschaft aller beobachteten
Prozesse: Die BABAR-Kollaboration berichtete eine direkte Beobachtung von
T-Verletzung im neutralen \(B\)-Meson-System \citep{BABAR2012}. Auch daraus
folgt jedoch weder Zeitreise noch ein steuerbarer Rückwärtskanal.

Mindestens vier Ebenen sind zu trennen:
\begin{enumerate}
  \item die Symmetrie des Bewegungsgesetzes;
  \item die Symmetrie oder Asymmetrie konkreter Rand- und Anfangsdaten;
  \item der thermodynamische Zeitpfeil makroskopischer Prozesse;
  \item die epistemische Richtung, in der Records erzeugt und ausgewertet
  werden.
\end{enumerate}
Ein mikroskopisch reversibles Gesetz kann mit makroskopisch irreversibler
Recordbildung vereinbar sein. Der thermodynamische Pfeil wird in statistischen
Beschreibungen wesentlich durch besondere Randbedingungen und die enorme
Asymmetrie zwischen typischen Mikro- und Makrozuständen bestimmt. Ein
gespeicherter Record ist ein physischer Zustand; sein Erzeugen, Kopieren oder
Löschen besitzt Trägerdynamik und kann thermodynamische Kosten haben
\citep{landauer1961}. Dass der Inhalt eines Records später anders bewertet
wird, ist davon zu unterscheiden. Modelle entgegengesetzter
thermodynamischer Zeitpfeile sind mathematisch untersuchbar
\citep{Schulman1999}; sie verlangen besondere Randbedingungen und sind kein
Nachweis, dass solche makroskopischen Pfeile in unserem beobachtbaren Bereich
tatsächlich zusammentreffen.

Auch eine formal unitäre Quantendynamik
\[
  \lvert\psi(t_2)\rangle
  =U(t_2,t_1)\lvert\psi(t_1)\rangle,
  \qquad
  U(t_1,t_2)=U(t_2,t_1)^{-1},
\]
ist invertierbar, solange das geschlossene System und der Hamiltonoperator
bekannt sind. Die rechnerische Rekonstruktion von
\(\lvert\psi(t_1)\rangle\) aus späteren Daten ist keine Nachricht, die bei
\(t_1\) empfangen wurde. Invertierbarkeit ist eine Eigenschaft der
mathematischen Dynamik, Retrokausalität eine Aussage über die Richtung
physischer Abhängigkeit.

\subsection{Retardierte und avancierte Lösungen klassischer Feldgleichungen}

Für eine lineare Feldgleichung schematischer Form
\begin{equation}
\label{eq:field-green}
  L\phi(x)=j(x)
\end{equation}
kann eine Lösung mittels einer Greenschen Funktion geschrieben werden:
\[
  \phi(x)=\phi_{\mathrm h}(x)
  +\int G(x,x')j(x')\,\dd x'.
\]
Eine retardierte Greensche Funktion \(G_{\mathrm{ret}}\) besitzt nur dann
Unterstützung, wenn \(x'\) in der kausalen Vergangenheit von \(x\) liegt. Eine
avancierte Greensche Funktion \(G_{\mathrm{adv}}\) verwendet demgegenüber
Quellenanteile aus der kausalen Zukunft. Beide können mathematisch dieselbe
lokale Differentialgleichung lösen; ihre physikalische Bedeutung wird durch
Randbedingungen und die Wahl der Lösungsklasse bestimmt.

Die Wheeler--Feynman-Absorbertheorie kombiniert retardierte und avancierte
Anteile in einer direkten Teilchenwechselwirkung
\citep{WheelerFeynman1945,WheelerFeynman1949}. Ihr historischer und
konzeptioneller Wert liegt gerade darin, dass eine zeit-symmetrische
Mikrodynamik und eine beobachtete Strahlungsasymmetrie nicht durch bloßes
Anschauen der lokalen Gleichung auseinandergehalten werden können. Benötigt
werden globale Absorberbedingungen und eine Analyse dessen, welche
beobachtbaren Freiheitsgrade verfügbar sind.

Die Existenz einer avancierten mathematischen Lösung beweist daher weder ihre
Realisierung in der Natur noch einen frei nutzbaren Rückwärtskanal. Umgekehrt
ist die bloße Ungewohntheit avancierten Verhaltens kein mathematischer
Widerspruch. Wissenschaftlich zu entscheiden ist, welche Randbedingungen die
Natur realisiert und ob daraus von der Standardbeschreibung unterscheidbare
Messwerte folgen.

\subsection{Vor- und Nachselektion in der Quantenmechanik}

In einem vor- und nachselektierten Experiment werde ein System zum Zeitpunkt
\(t_i\) im Zustand \(\lvert\psi_i\rangle\) präpariert und zum späteren
Zeitpunkt \(t_f\) auf \(\lvert\psi_f\rangle\) nachselektiert. Für eine ideale
Zwischenmessung mit Projektoren \(P_k\) ergibt die
Aharonov--Bergmann--Lebowitz-Regel
\citep{AharonovBergmannLebowitz1964}
\begin{equation}
\label{eq:abl}
 \mathbb P(a_k\mid\psi_i,\psi_f)
 =
 \frac{
 \left|\langle\psi_f\rvert U(t_f,t)P_kU(t,t_i)
 \lvert\psi_i\rangle\right|^2}
 {\displaystyle
 \sum_j
 \left|\langle\psi_f\rvert U(t_f,t)P_jU(t,t_i)
 \lvert\psi_i\rangle\right|^2}.
\end{equation}
Die spätere Auswahl \(\psi_f\) steht in dieser bedingten Wahrscheinlichkeit
explizit. Das ist mathematisch unstrittig. Es bedeutet zunächst, dass sich die
Statistik des Teilensembles mit festgelegtem Endergebnis von der Statistik des
Gesamtensembles unterscheidet.

Entsprechend kann für einen Operator \(A\) der schwache Wert
\begin{equation}
\label{eq:weak-value}
  A_w=
  \frac{\langle\psi_f\rvert A\lvert\psi_i\rangle}
       {\langle\psi_f\mid\psi_i\rangle}
\end{equation}
definiert werden. Schwache Werte können außerhalb des Eigenwertspektrums von
\(A\) liegen. Dies ist kein Beweis dafür, dass ein einzelnes System vor der
Nachselektion eine klassische, anomale Eigenschaft besessen habe. Die
physikalische Interpretation hängt vom Messmodell, von der Ensemblebildung
und vom gewählten ontologischen Rahmen ab
\citep{AharonovAlbertVaidman1988,Pusey2014}.

Der Zwei-Zustands-Vektor-Formalismus behandelt Vor- und Nachselektion
zeit-symmetrisch. Andere Interpretationen lesen dieselben bedingten
Wahrscheinlichkeiten ohne ontische Rückwirkung. Solange beide dieselben
beobachtbaren Verteilungen vorhersagen, entscheidet das Experiment nicht
zwischen ihnen. Eine Interpretation wird erst zu einer empirisch
ausgezeichneten Theorie, wenn sie eine unterscheidende Vorhersage liefert.

\subsection{Delayed Choice und Quantum Eraser}

Experimente mit verzögerter Wahl zeigen, dass die später gewählte
Messanordnung bestimmt, welche gemeinsame Statistik sinnvoll ausgewertet
wird. Beim Delayed-Choice-Quantum-Eraser werden frühere Detektionsereignisse
nach später verfügbaren, korrelierten Partnerergebnissen sortiert
\citep{ScullyDruehl1982,KimEtAl2000}. In den nachselektierten Teilmengen können Interferenzmuster
sichtbar werden, die in der unbedingten Summe verschwinden.
Realisierungen von Wheelers verzögerter Wahl und verzögert gewähltem
Entanglement Swapping bestätigen die kontextabhängigen
Korrelationsvorhersagen der Quantenmechanik
\citep{JacquesEtAl2007,MaEtAl2012}; sie stellen für sich keine frühere
lokale Nachricht bereit.

Seien \(R\) die früheren Detektordaten, \(B\) die spätere Basiswahl und \(Y\)
das spätere Partnerergebnis. Dann können durchaus
\[
  \mathbb P(R\mid B,Y)
  \neq
  \mathbb P(R)
\]
und unterschiedliche konditionierte Muster auftreten. Für einen
rückwärtsgerichteten Signalkanal wäre jedoch eine Abhängigkeit der bereits
lokal verfügbaren Marginalverteilung von der frei gewählten Basis erforderlich:
\begin{equation}
\label{eq:no-back-signal}
  \mathbb P(R\mid\operatorname{do}(B=b))
  =
  \mathbb P(R\mid\operatorname{do}(B=b'))
  =\mathbb P(R).
\end{equation}
In der Standardquantenmechanik bleibt diese No-signalling-Bedingung erhalten.
Das Muster entsteht erst durch den späteren Vergleich der Records und die
Auswahl korrelierter Teilensembles. Die früher gespeicherten Klicks werden
nicht überschrieben.

\statusbox{QAmber}{
\textbf{Interpretationsgrenze verzögerter Wahl.}
Delayed-Choice-Experimente schließen naive klassische Geschichten aus, in
denen ein Quantensystem unabhängig vom gesamten Messkontext stets eine
vorab festgelegte klassische Eigenschaft trägt. Sie beweisen jedoch keine
kontrollierbare Änderung eines früheren Records und keine Nachricht aus der
Zukunft.
}

\subsection{Bell-Korrelationen und retrokausale verborgene Variablen}

Für zwei räumlich getrennte Messstationen seien \(x,y\) die Einstellungen,
\(A,B\) die Ergebnisse und \(\lambda\) verborgene Variablen. Ein lokal
faktorisierendes Modell besitzt die Form
\begin{equation}
\label{eq:bell-factor}
 \mathbb P(A,B\mid x,y)
 =\int\dd\lambda\,
 \rho(\lambda\mid x,y)
 \mathbb P(A\mid x,\lambda)
 \mathbb P(B\mid y,\lambda).
\end{equation}
Die übliche Messunabhängigkeit setzt
\begin{equation}
\label{eq:measurement-independence}
  \rho(\lambda\mid x,y)=\rho(\lambda).
\end{equation}
Bell-Verletzungen zeigen, dass nicht alle gemeinsam verwendeten Annahmen
einer lokal-kausalen verborgenen-Variablen-Erklärung beibehalten werden
können \citep{Bell1964}. Retrokausale Modelle erlauben, dass spätere Einstellungen \(x,y\) an
der Bestimmung früherer oder zwischenliegender \(\lambda\)-Werte beteiligt
sind. Dadurch kann \cref{eq:measurement-independence} verletzt werden, ohne
eine unmittelbare räumliche Fernwirkung als fundamentalen Mechanismus zu
postulieren \citep{LeiferPusey2017,WhartonArgaman2020}.

Diese Strategie ist logisch und mathematisch ernst zu nehmen. Sie ist aber
nicht schon durch eine Bell-Verletzung bewiesen. Dieselben beobachtbaren
Korrelationen können in verschiedenen Interpretationen oder Ontologien
repräsentiert werden. Ein retrokausales Modell muss insbesondere erklären:
\begin{itemize}
  \item wie die Abhängigkeit von zukünftigen Einstellungen dynamisch oder
  randwerttheoretisch entsteht;
  \item warum sie Lorentz-kovariant ist;
  \item warum sie keine beobachtbare Rückwärtssignalisierung erzeugt;
  \item ob das erforderliche No-signalling robust oder nur durch feine
  Parameterabstimmung erhalten bleibt \citep{WoodSpekkens2015};
  \item welche neue Vorhersage das Modell gegenüber der Standardtheorie macht.
\end{itemize}

Retrokausalität und Superdeterminismus verletzen beide in möglichen
Darstellungen die statistische Unabhängigkeit von \(\lambda\) und den
Einstellungen. Ihr erklärter Mechanismus ist jedoch verschieden:
Superdeterministische Ansätze führen die Korrelation typischerweise auf
gemeinsame Bedingungen in der Vergangenheit zurück; retrokausale Ansätze auf
eine Abhängigkeit von späteren Messkontexten. Aus beobachtbaren
Wahrscheinlichkeiten allein ist diese ontologische Richtung nicht immer
identifizierbar. Deshalb muss sie durch Zusatzstruktur und, wenn möglich,
neue Messfolgen operationalisiert werden.

\subsection{Interpretations- und Modellfamilien}

Die Literatur enthält mehrere Familien, in denen zeit-symmetrische oder
retrokausale Begriffe eine tragende Rolle spielen. Sie sollten nicht unter
einem einzigen Wahrheitsstatus zusammengefasst werden.

\begin{longtable}{>{\raggedright\arraybackslash}p{0.20\textwidth}
                  >{\raggedright\arraybackslash}p{0.25\textwidth}
                  >{\raggedright\arraybackslash}p{0.22\textwidth}
                  >{\raggedright\arraybackslash}p{0.22\textwidth}}
\toprule
\textbf{Familie} & \textbf{Zeitliche Struktur} & \textbf{Rückwärtssignal?} & \textbf{Epistemischer Status} \\
\midrule
\endfirsthead
\toprule
\textbf{Familie} & \textbf{Zeitliche Struktur} & \textbf{Rückwärtssignal?} & \textbf{Epistemischer Status} \\
\midrule
\endhead
Vor-/Nachselektion und TSVF & Anfangs- und Endzustand konditionieren Zwischenstatistik & nein in der Standardanwendung & konsistenter Formalismus; Ontologie interpretations\-abhängig \\
Transaktionale Interpretation & Angebot- und Bestätigungsanteile bilden eine zeit-symmetrische Transaktion & kein frei nutzbarer Kanal vorgesehen & Interpretation der Quantenmechanik \citep{Cramer1986} \\
retrokausale Hidden-Variable-Modelle & spätere Einstellungen beeinflussen verborgene Zwischen- oder Frühvariablen & typischerweise durch No-signalling ausgeschlossen & aktive Modellforschung; nicht empirisch eindeutig ausgewählt \\
globale Raumzeit- oder Randwertmodelle & ganze Geschichte erfüllt lokale Gleichungen und mehrere Randbedingungen & nicht notwendig & mathematisch mögliche Repräsentation; Kausalinterpretation zusätzlich \\
retrokausale Feldmodelle & Wahrscheinlichkeitsmaß über Raumzeitfelder mit zeitlich beidseitiger Struktur & modellabhängig & konkrete aktive Forschungsmodelle; keine allgemeine Bestätigung \citep{WhartonArgaman2020} \\
Absorbertheorie & avancierte und retardierte direkte Wechselwirkung & makroskopische Signalfrage durch globale Bedingungen beschränkt & historisch bedeutendes zeit-symmetrisches Modell \\
\bottomrule
\end{longtable}

Aus dieser Vielfalt folgt ein methodischer Grundsatz: Der Satz ``Die
Quantenmechanik ist retrokausal'' ist ohne Angabe des Formalismus, der
ontologischen Variablen, der Randbedingungen und der operationalen
Konsequenzen unterbestimmt. Wissenschaftlich prüfbar ist nur eine konkrete
Modellklasse mit einer expliziten Wahrscheinlichkeits- oder
Bewegungsstruktur.

\subsection{Indefinite Kausalordnung ist nicht dasselbe wie Retrokausalität}

In Prozessmatrix- und Quantum-Switch-Formulierungen kann die globale Ordnung
zweier Operationen unbestimmt oder kohärent kontrolliert sein
\citep{OreshkovCostaBrukner2012}. Das bedeutet, dass der Gesamtprozess nicht notwendig als
klassische Mischung einer einzigen festen Reihenfolge dargestellt werden
kann. Daraus folgt nicht automatisch, dass ein Ereignis eine Nachricht an
seine eigene Vergangenheit sendet.

Zu unterscheiden sind:
\begin{enumerate}
  \item keine feste Reihenfolge zwischen zwei lokalen Operationen;
  \item eine Überlagerung oder prozessuale Unbestimmtheit von Ordnungen;
  \item ein geschlossener gerichteter Kausalpfad;
  \item ein operationaler Kanal zu einem bereits früher auslesbaren Record.
\end{enumerate}
Nur die letzten beiden Punkte berühren unmittelbar kausale Schleifen oder
Rückwärts\-signalisierung. Der Begriff ``indefinit'' darf nicht als Synonym für
``rückwärts'' verwendet werden.

\subsection{Geschlossene zeitartige Kurven und Konsistenzprobleme}

Allgemein-relativistische Lösungen können unter besonderen globalen
Voraussetzungen geschlossene zeitartige Kurven enthalten, beispielsweise in
Gödels rotierendem Kosmos \citep{Goedel1949}; traversierbare
Wurmloch-Konstruktionen machen zusätzlich die Rolle von Energiebedingungen
sichtbar \citep{MorrisThorneYurtsever1988}. Ein idealisierter
Beobachter könnte dann entlang einer stets lokal zukunftsgerichteten Weltlinie
in den eigenen kausalen Vergangenheitsbereich gelangen. Das ist ein anderer
Mechanismus als eine avancierte Feldlösung oder Nachselektion.

Quantentheoretische Modelle geschlossener zeitartiger Kurven verwenden unter
anderem Selbstkonsistenzbedingungen für eine Dichtematrix
\citep{Deutsch1991}. Schematisch muss ein Zustand \(\rho_C\) auf dem
geschlossenen Anteil eine Fixpunktbedingung erfüllen:
\begin{equation}
\label{eq:ctc-fix}
  \rho_C=\operatorname{Tr}_S\!\left[
  U(\rho_S\otimes\rho_C)U^\dagger
  \right].
\end{equation}
Eine solche Bedingung kann logische Paradoxien vermeiden, verändert aber die
gewöhnliche Theorie offener Anfangswertprobleme. Andere Vorschläge verwenden
Postselektion, Selbstkonsistenz oder verzweigende Geschichten. Hawking
formulierte demgegenüber die Chronology-Protection-Vermutung, nach der
Quanteneffekte die Ausbildung makroskopisch nutzbarer geschlossener
zeitartiger Kurven verhindern könnten \citep{Hawking1992}.

Weder die mathematische Existenz einer Raumzeitlösung noch die Konsistenz
eines idealisierten Quantenmodells belegt, dass solche Kurven in der Natur
erzeugbar sind. Dafür wären Energiebedingungen, Stabilität, Rückwirkung,
Quanteneffekte und eine physikalische Herstellungssituation gemeinsam zu
untersuchen.

\subsection{Paradoxien, Konsistenz und Informationsbegriffe}

Das Großvaterparadox entsteht, wenn eine frei wählbare spätere Handlung genau
die Bedingungen ihrer eigenen Möglichkeit beseitigt. Das Bootstrap-Paradox
entsteht, wenn Information oder ein Objekt in einer Schleife vorkommt, ohne
einen unabhängigen Ursprung zu besitzen. Diese Paradoxien zeigen, dass eine
naive Kombination von beliebiger Intervention, eindeutiger Geschichte und
Rückwärtskanal inkonsistent sein kann. Sie beweisen nicht, dass jede
retrokausale Modellierung unmöglich ist.

Mögliche mathematische Antworten sind:
\begin{itemize}
  \item \emph{Selbstkonsistenz}: Nur global widerspruchsfreie Geschichten sind
  zulässig;
  \item \emph{eingeschränkte Interventionen}: Nicht jede lokal vorstellbare
  Handlung ist mit den Randbedingungen kompatibel;
  \item \emph{Verzweigung}: Die Wirkung trifft nicht dieselbe eindeutige
  Geschichte;
  \item \emph{probabilistische Fixpunkte}: Konsistenz gilt für Verteilungen,
  nicht für jede Einzelgeschichte;
  \item \emph{No-signalling}: Eine ontische Rückabhängigkeit kann bestehen,
  ohne als Nachricht kontrollierbar zu sein.
\end{itemize}

Ein semantischer Informationsgehalt ist zudem vom physischen Signal zu
trennen. Für einen Rückwärtskanal genügt es nicht, dass ein heutiger Beobachter
eine frühere Spur im Licht späterer Daten neu versteht. Erforderlich wäre ein
physischer Trägerzustand bei \(t_1\), dessen auslesbare Statistik von einer bei
\(t_2\) frei vorgenommenen Codierung abhängt. Ohne Träger, Kopplung und
Decodierverfahren liegt keine operative Nachrichtenübertragung vor.

\subsection{Der fehlende Rückwärtskanal im gegenwärtigen Mandelbrot-Prozess}

Für das in \cref{sec:dynamische-synthese} definierte Modell gilt
\[
  z_0(c)=0,
  \qquad
  z_{n+1}(c)=z_n(c)^2+c.
\]
Der endliche Gate-Status \(\mathcal A_N(c)\) wird aus dem Parameter und dem
bis \(N\) berechneten Präfix bestimmt. Er ist kein Eingang der Rekursion.

\begin{satz}[Kausale Autonomie früher Präfixe im definierten Prozess]
\label{satz:kein-rueckwaertskanal}
Seien \(c\in U\) und \(N\in\N_0\) fest. Für jedes \(k\le N\) ist
\(z_k(c)\) ausschließlich durch \(c\), \(z_0=0\) und die \(k\)-fache
Anwendung von \(z\mapsto z^2+c\) bestimmt. Änderungen der späteren
Auswertungsstufe, der Gate-Abfrage oder eines Status
\(\mathcal A_M(c)\) mit \(M>k\) verändern \(z_k(c)\) nicht. Das gegenwärtig
definierte Mandelbrot-Prozessmodell enthält daher keinen gerichteten Kanal von
einem späteren Gate-Ausgang zu einem früheren Rekursionszustand.
\end{satz}

\begin{proof}
Der Beweis erfolgt durch Induktion über \(k\). Für \(k=0\) gilt
\(z_0(c)=0\) unabhängig von jeder Gate-Auswertung. Sei die Behauptung für
\(k\) gezeigt. Dann ist
\[
  z_{k+1}(c)=z_k(c)^2+c
\]
allein eine Funktion des bereits bestimmten Zustands \(z_k(c)\) und des
festen Parameters \(c\). In der rechten Seite tritt weder
\(\mathcal A_M(c)\) noch eine spätere Auswertungsentscheidung auf. Folglich
ist auch \(z_{k+1}(c)\) davon unabhängig. Induktiv gilt dies für alle
\(k\le N\).

Der Gate-Status ist umgekehrt eine Funktion vorhandener Zustände und
gegebenenfalls zusätzlicher Zertifikate. Der funktionale Pfeil verläuft somit
im definierten Modell von der Spur zum Status, nicht vom späteren Status zu
früheren Spurwerten. Eine Umkehr dieses Pfeils müsste als zusätzliche
Dynamik ausdrücklich definiert werden.
\end{proof}

\begin{korollar}[Reklassifikation ohne Überschreiben]
\label{kor:reklassifikation-ohne-ueberschreiben}
Ein Wechsel von \(\CONTINUE\) zu \(\PASS\) oder \(\BLOCK\) kann die
zulässige Aussage über das gesamte Präfix ändern, aber keinen in diesem
Präfix enthaltenen Wert überschreiben. \Cref{prop:rueckbestimmung} ist damit
nicht nur eine sprachliche Unterscheidung, sondern eine Folge der gerichteten
Funktionsabhängigkeit des Modells.
\end{korollar}

Sei zusätzlich \(E_M\) eine frei wählbare spätere Entscheidung, ob und wie
weit das Gate ausgewertet wird, wobei \(E_M\) weder \(c\) noch die Rekursion
\(F\) verändert. Für jeden früher erzeugten Record
\(R_k=r(c,z_0,\ldots,z_k)\), \(k<M\), gilt dann modellintern
\begin{equation}
\label{eq:qik-no-back-signal}
 \mathbb P(R_k\mid\operatorname{do}(E_M=e))
 =
 \mathbb P(R_k\mid\operatorname{do}(E_M=e')).
\end{equation}
Diese Gleichung ist eine strukturelle Konsequenz der aktuellen
Modellarchitektur. Sie ist kein allgemeines Naturgesetz. Wenn \(E_M\) den
Parameter \(c\), die Anfangsbedingung oder eine physische Kopplung verändert,
wird ein anderes Modell betrachtet.

Eine tatsächlich retrokausale Erweiterung müsste mindestens eine der
folgenden Strukturen ergänzen:
\begin{enumerate}
  \item eine Endrandbedingung \(B_f\), von der frühere Zustände abhängen;
  \item einen Übergangskern
  \(K(z_k\mid z_{k-1},c,A_{t_2})\) mit expliziter späterer Einstellung;
  \item eine globale Konsistenzbedingung für vollständige Geschichten;
  \item eine physikalische Korrespondenz \(\Phi\), in der spätere
  Raumzeitbedingungen einen früheren Trägerzustand beeinflussen.
\end{enumerate}
Keine dieser Strukturen folgt aus der quadratischen Iteration oder aus der
Shift-Invarianz der Klassifikation allein.

\statusbox{QGreen}{
\textbf{Bewiesen für das definierte Prozessmodell.}
Die spätere Gate-Auswertung verändert keinen früheren Rekursionszustand. Das
Modell realisiert semantische Rückbestimmung, aber keinen Zustandskanal aus
einem späteren Gate in ein früheres Präfix.
}

\subsection{Wie eine physikalische Retrokausalitätshypothese formuliert werden müsste}

Eine physikalische Erweiterung darf nicht allein aus einer zeitlich
rückwärts lesbaren Grafik bestehen. Sie benötigt einen Zustandsraum
\(\mathcal Z\), eine Dynamik oder ein globales Variationsprinzip, Randdaten,
Observablen und eine Wahrscheinlichkeitsstruktur. Eine schematische
Zwei-Randwertformulierung kann lauten
\begin{equation}
\label{eq:two-boundary}
  \delta S[q]=0,
  \qquad
  q(t_i)=q_i,
  \qquad
  B_f(q(t_f),A_{t_f})=0.
\end{equation}
Die Lösungen \(q(t)\) hängen dann mathematisch von beiden Randbedingungen ab.
Ob \(A_{t_f}\) als frei kontrollierbare zukünftige Ursache gelten darf, ist
eine zusätzliche physikalische und kausaltheoretische Frage.

Für eine statistische Feldtheorie könnte entsprechend ein Maß über ganze
Geschichten verwendet werden:
\begin{equation}
\label{eq:path-retro}
  \mathbb P[\phi\mid B_i,B_f]
  =\frac{1}{Z[B_i,B_f]}
  \exp\!\left(-\frac{1}{\hbar}\mathcal I[\phi]\right)
  \mathbf 1_{B_i}[\phi]\mathbf 1_{B_f}[\phi],
\end{equation}
wobei \(\mathcal I\) hier nur symbolisch für eine geeignete reelle oder
euklidische Wirkung steht. Ein physikalisches Modell muss Vorzeichen,
Konturen, Eichsymmetrien, Normierbarkeit und Lorentz-Signatur konkret
angeben. Eine Formel mit einem unbestimmten Funktional ist noch keine Theorie.

Für die Verbindung zum dimensionslosen Mandelbrot-Modell müsste zusätzlich
eine Abbildung
\[
  \Phi_{B_i,B_f}:\mathcal Z_{\mathrm{phys}}\longrightarrow U
\]
definiert werden. Sie müsste Dimensionen, Symmetrien, Kausalstruktur und
Observablen erhalten oder ihre kontrollierte Abweichung quantifizieren. Erst
dann wäre prüfbar, ob ein retrokausales physikalisches Modell tatsächlich an
die QIK--VRT-Anschlussordnung semikonjugiert ist.

\subsection{Experimentelles Mindestprotokoll für einen Rückwärtskanal}
\label{sec:retro-protokoll}

Ein überzeugender Test muss eine spätere freie Einstellung von einem früheren
Record trennen und jeden gewöhnlichen Informationspfad schließen. Ein
Minimalprotokoll umfasst folgende Schritte.

\begin{enumerate}
  \item \textbf{Präparation bei \(t_0\).}
  Der Anfangszustand, alle Kalibrierparameter und zulässigen Ausschlussregeln
  werden vorab dokumentiert.

  \item \textbf{Frühe Messung bei \(t_1\).}
  Ein Rohrecord \(R_{t_1}\) wird erzeugt. Sein vollständiger Datenstrom wird
  unverändert gespeichert, gehasht, extern zeitgestempelt und vor der späteren
  Wahl versiegelt.

  \item \textbf{Spätere randomisierte Intervention bei \(t_2>t_1\).}
  Die Einstellung \(A_{t_2}\) wird erst nach der Versiegelung erzeugt. Die
  Quelle der Wahl, ihre statistischen Eigenschaften und ihre möglichen
  Seitenkanäle werden offengelegt. Die Randomisierung muss unabhängig von den
  Personen erfolgen, die die frühen Daten auswerten.

  \item \textbf{Späte Messung bei \(t_3\).}
  Späte Daten \(Y_{t_3}\) werden getrennt erfasst. Postselektion ist entweder
  ausgeschlossen oder vollständig vorregistriert; die unbedingte frühe
  Marginalverteilung bleibt der primäre Endpunkt.

  \item \textbf{Blindes Öffnen bei \(t_4\).}
  Erst nach Abschluss aller Läufe werden die Zuordnungen zwischen
  \(R_{t_1}\) und \(A_{t_2}\) geöffnet. Analysecode, Ausschlusskriterien,
  Teststatistik und Stichprobengröße müssen zuvor fixiert sein.
\end{enumerate}

Die Nullhypothese lautet für alle Einstellungen \(a,a'\)
\begin{equation}
\label{eq:retro-null}
 H_0:\quad
 \mathbb P(R_{t_1}\mid\operatorname{do}(A_{t_2}=a))
 =
 \mathbb P(R_{t_1}\mid\operatorname{do}(A_{t_2}=a')).
\end{equation}
Eine Alternative \(H_R\) muss vorab angeben, welche Recordkomponente sich um
welchen Betrag und mit welcher Abhängigkeit von \(a\) verschiebt. Eine bloße
nachträgliche Suche nach irgendeiner abweichenden Untergruppe ist kein
belastbarer Test.

Für diskrete Records kann beispielsweise die totale Variation
\begin{equation}
\label{eq:retro-tv}
  \Delta_{a,a'}
  :=\frac12\sum_r
  \left|
  \mathbb P(r\mid\operatorname{do}(a))
  -\mathbb P(r\mid\operatorname{do}(a'))
  \right|
\end{equation}
als Effektgröße dienen. Rückwärtssignalisierung würde
\(\Delta_{a,a'}>0\) erfordern, nachdem alle vorwärts gerichteten Kanäle und
systematischen Verzerrungen berücksichtigt wurden. Neben einem Signifikanztest
sind Konfidenz- oder glaubwürdige Intervalle, vorher festgelegte
Mindest-Effektgrößen und Replikationen an unabhängigen Orten erforderlich.

Mindestens folgende Fehlerquellen müssen durch Kontrollmessungen oder
Architektur ausgeschlossen werden:
\begin{itemize}
  \item Clock- und Synchronisationsfehler;
  \item Cache-, Netzwerk-, Datei- und Metadatenleckage;
  \item unbemerkte Abhängigkeit der Randomisierung vom frühen System;
  \item gemeinsame Umweltursachen und Drift;
  \item selektiver Verlust von Läufen;
  \item flexible Wahl von Zeitfenstern, Filtern oder Untergruppen;
  \item nachträgliche Änderung des Analyseprogramms;
  \item gewöhnliche Quantenkorrelationen, die erst nach klassischem
  Datenvergleich sichtbar werden;
  \item Manipulation oder Überschreiben des frühen Records nach \(t_1\).
\end{itemize}

Ein positiver Einzelbefund wäre zunächst ein Anomaliebericht. Erst die
Wiederholung mit neu erzeugten Einstellungen, unabhängiger Hardware,
unabhängigen Auswertenden und offengelegtem Code könnte eine retrokausale
Interpretation gegenüber systematischen Fehlern stärken. Zudem müsste geprüft
werden, ob eine vorwärtskausale Erweiterung des Modells dieselben Daten
erklärt.

\subsection{Was eine Bestätigung ändern würde --- und was nicht}

Würde die Bedingung \eqref{eq:retro-ontisch} für eine klar definierte ontische Variable
empirisch ausgezeichnet, hätte dies weitreichende Folgen für physikalische
Erklärungen. Anfangswertformulierungen müssten möglicherweise durch globale
Randwert- oder Konsistenzstrukturen ergänzt werden. Kausale Modelle dürften
nicht mehr voraussetzen, dass jede gerichtete Abhängigkeit der
Koordinatenzeit folgt. In Quantenfundamenten könnte eine lokal vermittelte,
zeit-symmetrische Erklärung von Bell-Korrelationen an Gewicht gewinnen.

Eine solche Bestätigung würde jedoch nicht automatisch bedeuten:
\begin{itemize}
  \item dass Nachrichten in die Vergangenheit gesendet werden können;
  \item dass historische Records beliebig umgeschrieben werden;
  \item dass Energie- oder Impulserhaltung aufgehoben ist;
  \item dass logische Widersprüche physisch realisiert werden;
  \item dass menschliche Entscheidungen bedeutungslos oder vollständig
  vorbestimmt sind;
  \item dass jedes zeit-symmetrische oder fraktale Modell wahr ist;
  \item dass QIK--VRT ohne eine neue Korrespondenzabbildung eine bestätigte
  Naturtheorie wäre.
\end{itemize}

Würde sogar die Bedingung \eqref{eq:retro-signal} reproduzierbar erfüllt, wären die Folgen
noch stärker. Grundannahmen über kausale Kommunikation, zeitgebundene
kryptografische Zusagen und die Zuordnung von Verantwortung müssten neu
untersucht werden. Zugleich wären Konsistenzbedingungen unverzichtbar: Ein
solcher Kanal ist nur physikalisch beschreibbar, wenn die Gesamtdynamik für
alle tatsächlich zulässigen Eingaben widerspruchsfreie Wahrscheinlichkeiten
erzeugt.

Das Ausbleiben eines solchen Signals widerlegt hingegen nicht jede
retrokausale Ontologie. Wenn eine Modellklasse konstruktionsbedingt exakt
No-signalling erfüllt, sind zusätzliche beobachtbare Kriterien nötig. Ohne
unterscheidende Vorhersage bleibt die Wahl zwischen ontologisch verschiedenen,
aber empirisch äquivalenten Darstellungen unterbestimmt.

\subsection{Konsequenz für die Ontologie des Unterschieds}

Die Ontologie des universalen Unterschieds gewinnt an wissenschaftlicher
Schärfe, wenn sie diese Ebenen nicht zusammenführt. Der relevante Unterschied
ist nicht pauschal ``Zukunft gegen Vergangenheit'', sondern mindestens
\[
\begin{gathered}
  \text{früherer Zustand}
  \neq\text{früherer Record}
  \neq\text{spätere Evidenz},\\
  \text{Korrelation}
  \neq\text{Intervention}
  \neq\text{Signal},\\
  \text{globale Konsistenz}
  \neq\text{lokaler Wirkmechanismus},\\
  \text{interpretative Möglichkeit}
  \neq\text{empirische Bestätigung}.
\end{gathered}
\]

``Transzendierende Information'' kann im gegenwärtigen QIK--VRT-Modell
präzise bedeuten, dass ein Record erhalten bleibt, spätere Evidenz mit ihm
verbunden wird und die zulässige Klassifikation der gesamten Geschichte
stabilisiert. Diese Wirkung ist technisch real: Sie verändert Entscheidungen,
Freigaben und Folgewirkungen nach dem Zeitpunkt der neuen Evidenz. Sie wirkt
aber nicht dadurch, dass der frühere physische Record rückwirkend
überschrieben wird.

Gerade diese Selbstbegrenzung ist keine Schwächung, sondern ein
erkenntnistheoretischer Gewinn. Sie stellt sicher, dass ein semantisch
reichhaltiger Prozess nicht allein aufgrund seiner sprachlichen Form als neue
physikalische Kraft ausgegeben wird. Zugleich lässt sie den Raum für eine
konkrete retrokausale Hypothese offen, sofern deren Zustände, Kopplungen,
Einheiten, Randbedingungen und Messfolgen vollständig angegeben werden.

\subsection{Statusbilanz}

\statusbox{QGreen}{
\textbf{DONE.}
Begrifflich getrennt sind semantische Rückbestimmung, Retrodiktion,
Nachselektion, Zeitumkehrsymmetrie, avancierte Lösung, ontische
Retrokausalität, geschlossene zeitartige Kurve und operative
Rückwärtssignalisierung. Für den gegenwärtig definierten
Mandelbrot-/QIK--VRT-Prozess ist bewiesen, dass spätere Gate-Ausgänge frühere
Rekursionszustände nicht verändern. Delayed-Choice-Korrelationen und
vor-/nachselektierte Wahrscheinlichkeiten sind ohne einen Rückwärtskanal
mathematisch darstellbar.
}

\statusbox{QAmber}{
\textbf{CONTINUE.}
Zeit-symmetrische, Zwei-Randwert- und retrokausale Modelle sind
wissenschaftlich diskutierbare Modellklassen. Offen sind eine konkrete
dimensionsrichtige QIK--VRT-Erweiterung, ihre Lorentz-kovariante Dynamik, die
Erhaltung von No-signalling oder eine begründete Abweichung davon, eine
unterscheidende Vorhersage sowie unabhängige experimentelle Reproduktion.
}

\statusbox{QRed}{
\textbf{BLOCK.}
Nicht zulässig sind die Schlüsse von späterer Reklassifikation auf eine
Änderung früherer Zustände, von Bayes-Konditionierung auf physische
Rückwirkung, von Zeitumkehrinvarianz auf eine zukünftige Ursache, von
Delayed-Choice- oder Quantum-Eraser-Daten auf einen nutzbaren Rückwärtskanal,
von indefiniter Kausalordnung auf Zeitreise oder von der mathematischen
Konsistenz eines retrokausalen Modells auf dessen empirische Wahrheit.
}

Die wissenschaftlich vollständige Schlussform dieses Abschnitts lautet:
\[
\boxed{
\begin{aligned}
\text{QIK--VRT-Reklassifikation}
&\Longrightarrow
\text{semantische Rückbestimmung},\\
\text{QIK--VRT-Reklassifikation}
&\not\Longrightarrow
\text{physikalische Retrokausalität},\\
\text{physikalische Retrokausalität}
&\not\Longrightarrow
\text{Rückwärtssignalisierung},\\
\text{Rückwärtssignalisierung}
&\Longrightarrow
\text{neue, unabhängig zu reproduzierende Physik}.
\end{aligned}}
\]

\clearpage
\part{Geltungsprüfung, Ontologie und Forschungsprogramm}

\section{Falsifizierbare physikalische Mindestanforderungen}
\label{sec:phys-falsifikation}

Eine physikalische QIK--VRT- oder Mandelbrot-Raumzeit-Hypothese
\(H_{\mathrm Q}\) wird gegenüber einer Referenztheorie \(H_0\) erst
empirisch unterscheidbar, wenn für mindestens eine vorab bestimmte Observable
\(O\) und zugehörige Daten \(D\) gilt
\[
  P(D\mid H_{\mathrm Q})\neq P(D\mid H_0)
\]
in einem Bereich, in dem die Differenz nicht durch Messunsicherheit, freie
Nachanpassung von \(\Phi\) oder bekannte systematische Effekte absorbiert
werden kann.

Dafür sind mindestens erforderlich:
\begin{enumerate}
  \item eine vollständig angegebene, dimensionsrichtige und
  koordinatenunabhängige Korrespondenz \(\Phi\);
  \item eine definierte Dynamik \(G_{\Delta\tau}\) und ein präziser
  Zusammenhang mit \(F\);
  \item vor Betrachtung der Testdaten festgelegte Referenzskalen, Parameter
  und Toleranzen;
  \item die Rückgewinnung etablierter Grenzfälle von Allgemeiner
  Relativitätstheorie und Quantenfeldtheorie;
  \item eine Vorhersage mit Einheit, Zahlenwert, Unsicherheit und einer von
  \(H_0\) unterscheidbaren Abhängigkeit;
  \item numerische Konvergenz unter Verfeinerung von Raumgitter, Zeitschritt
  und Iterationsgrenze;
  \item unabhängige Reproduktion von Daten, Code, Parametern und Auswertung.
\end{enumerate}

Eine nachträglich frei wählbare Abbildung \(\Phi\), die jedes beobachtete
Muster auf irgendeinen Fraktalausschnitt abbilden kann, wäre nicht
falsifizierbar. Wissenschaftliche Anschlussfähigkeit verlangt deshalb nicht
nur Ausdrucksstärke, sondern Beschränkung.

\statusbox{QGreen}{
\textbf{DONE.}
Dimensionshomogene Normierung, Unterscheidung zwischen positiver
Darstellungsmetrik und Lorentz-Metrik, SI-Dimensionen der
Einstein-Gleichungen, Planck-Einheiten sowie Schwarzloch-Entropie und
Hawking-Temperatur.
}

\statusbox{QAmber}{
\textbf{CONTINUE.}
Konstruktion einer invarianztreuen Abbildung \(\Phi\), Nachweis einer
Semikonjugation, gemeinsame Testvektoren, numerische Konvergenz und Ableitung
einer unterscheidenden physikalischen Vorhersage.
}

\statusbox{QRed}{
\textbf{BLOCK.}
Der Schluss von Dimensionsrichtigkeit auf Kausalität, vom Rechengitter auf
ontische Raumzeitquanten, von visueller Selbstähnlichkeit auf Naturidentität,
von einem blockierten Update auf einen physikalischen Fixpunkt oder von
späterer Reklassifikation auf rückwärts gerichtete Signalkausalität.
}

Der entscheidende Unterschied dieses dritten Teils lautet
\[
\boxed{
\begin{aligned}
\text{Einheit}&\neq\text{Größe},\\
\text{Dimensionsgleichung}&\neq\text{Naturgesetz},\\
\text{Lorentz-Metrik}&\neq\text{positiver Darstellungsabstand},\\
\text{numerischer Zeitschritt}&\neq\text{Zeitquant},\\
\text{Information}&\neq\text{Energie ihres Trägers},\\
\text{Modellkorrespondenz}&\neq\text{empirische Identität}.
\end{aligned}}
\]

\section{Die Ontologie des universalen Unterschieds}
\label{sec:ontologie}

\subsection{Transzendentaler Kern}

Eine bestimmte Aussage muss von ihrer Verneinung, ein Term von einem anderen
Term und ein Beweis von einem Nichtbeweis unterscheidbar sein. Auch die
Bestreitung des Unterschieds muss sich von ihrer Bejahung unterscheiden, um
überhaupt etwas Bestimmtes zu behaupten. Die folgende Folge wird als
transzendentales Argument formuliert, nicht als zusätzliches Theorem der
Mengenlehre oder der empirischen Naturwissenschaft. In diesem Sinn gilt
\[
  \text{Bestimmtheit}
  \Longrightarrow
  \text{Unterscheidbarkeit}
  \Longrightarrow
  \text{Unterschied}.
\]

Der mengentheoretische Komplementsatz liefert dazu eine präzise formale
Instanz. Sobald ein Grundbereich \(U\) und eine bestimmte Klasse
\(A\subseteq U\) festgelegt sind, ist zugleich die vollständige
Unterscheidung
\[
  U=A\dotcup(U\setminus A)
\]
bestimmt. Das Ganze und der Unterschied schließen einander nicht aus. Das
Ganze wird vielmehr als vollständige Vermittlung komplementärer Bestimmtheit
darstellbar.

\subsection{Der entscheidende Unterschied}

Für den hier untersuchten Fall lautet die zentrale ontologische Anwendung:
\[
\begin{aligned}
\text{Menge} &\neq \text{Komplement},\\
\text{Teilbereiche} &\neq \text{Ganzes},\\
\text{Modelluniversum} &\neq \text{empirischer Kosmos},\\
\text{Quantisierbarkeit} &\neq \text{Quantisierung},\\
\text{Darstellbarkeit} &\neq \text{Beobachtbarkeit},\\
\text{Archivierung} &\neq \text{Validierung}.
\end{aligned}
\]

Diese Unterschiede begrenzen die These nicht von außen. Sie sind ihre
Selbstanwendung: Die Ontologie bleibt nur dann konsistent, wenn sie auch die
Geltungsarten ihrer eigenen Aussagen unterscheidet.

\subsection{Universales Schöpfungsprinzip}

Als ontologische Interpretation kann folgende Struktur formuliert werden:
\[
  \text{Unterschied}
  \longrightarrow
  \text{Bestimmtheit}
  \longrightarrow
  \text{Relation}
  \longrightarrow
  \text{Iteration}
  \longrightarrow
  \text{neue Struktur}.
\]

Die Pfeile bezeichnen hier ein strukturelles Interpretationsschema und keine
aus dem Komplementsatz allein deduzierte Kausalkette. Komplementbildung und
quadratische Iteration liefern eine formal präzise Instanz dieses Schemas;
sie beweisen weder dessen ausschließliche Universalität noch eine physikalische
Kosmogenese.

\begin{definition}[Universales Schöpfungsprinzip, strukturelle Lesart]
Das \emph{universale Schöpfungsprinzip} bezeichnet die These, dass bestimmte
neue Strukturen durch gesetzte oder entstandene Unterschiede, ihre Erhaltung,
ihre Relationierung und ihre rekursive Fortsetzung hervorgebracht werden.
\end{definition}

Diese Definition ist in Gegenstandsbereichen mit ausdrücklich definierten
Zuständen, Unterschieden, Relationen und Fortsetzungsregeln formal
instanziierbar. Sie bezeichnet keine zusätzlich postulierte Naturkraft. Der
formale Komplementsatz und die Iterationsstruktur machen die Deutung
mathematisch präzise. Dass sie zugleich die vollständige empirische Kosmogenese
beschreibt, wäre eine stärkere physikalische Hypothese und verlangte die in
\cref{sec:forschungsprogramm} angegebenen Tests.

\subsection{Der zweite Teilaspekt: Unterschied als geprüfte Fortsetzung}

Der dynamische Erkenntnisgewinn geht über die statische Formel
\(A\dotcup(U\setminus A)=U\) hinaus. Ein Unterschied wird im Prozess nicht
nur bezeichnet, sondern als mögliche Wirkung fortgesetzt, geprüft und in
einen revidierbaren Status überführt:
\[
\boxed{
\text{erste unterscheidbare Spur}
\longrightarrow F^n
\longrightarrow\text{Kandidat}
\longrightarrow\mathcal A_N
\longrightarrow\text{Status}
\longrightarrow\text{Rückkopplung}.}
\]
Im Mandelbrot-Demonstrator ist der Übergang \(0\to1\) exakt als erstmaliges
Auftreten eines Fluchtzertifikats darstellbar. Das Zertifikat trennt eine
bislang offene Spur von einer bewiesenen Außenklassifikation. Die
Gesamtordnung stabilisiert sich nicht notwendig, weil jeder Einzelzustand
konvergiert, sondern weil die richtige Endklassifikation gegenüber dem
Entfernen endlich vieler Anfangsschritte invariant ist:
\[
  \widehat\chi\circ\sigma=\widehat\chi.
\]

Damit erhält die Kurzform
\[
  \Omega=\Fix(\mathcal S)
\]
einen wissenschaftlich bestimmten Sinn, wenn \(\mathcal S\) den geprüften
Klassifikationsoperator bezeichnet und \(\Omega\) das in
\cref{sec:dynamische-synthese} definierte Prozessmodelluniversum und
\(\Fix\) den Fixpunkt des angegebenen Klassifikationsoperators bezeichnet.
Sie darf nicht ohne weitere Definition als Konvergenz jedes Orbits, als
Nash-Gleichgewicht oder als fertige physikalische Weltgleichung gelesen
werden.

Auch die Rede von ``transzendierender Information'' erhält so eine präzise
Lesart: Ein späteres Zertifikat kann die zulässige Bedeutung einer früheren
Spur ändern, während deren gespeicherte Zustände unverändert bleiben. Die
Information wirkt über den Augenblick hinaus, weil sie erhalten, verbunden
und erneut ausgewertet wird. Das ist echte Prozesswirkung im
Informationssystem, aber noch keine Energie- oder Signalübertragung in die
Vergangenheit. Der Unterschied zwischen semantischer Rückbestimmung und
physikalischer Retrokausalität ist die Selbstbegrenzung der These.

\section{Öffentliche Evidenz und Vorarbeiten}
\label{sec:evidenz}

Die hier entwickelte These steht in einer öffentlich prüfbaren
Artefaktkette. Die frühe Zenodo-Version von \qikvrt{} ist als Software unter
dem versionsspezifischen DOI \doilink{10.5281/zenodo.20712301}
archiviert \citep{lohmannQIKVRT833}. Die einschlägigen mathematischen und
quanten\-gravitativen Ausarbeitungen befinden sich als Dateien innerhalb des
Zenodo-Software-Releases \emph{RFC-QIKVRT-0001}, Version
\texttt{v2.13.4av1-node-open-multi-node}, unter dem versionsspezifischen DOI
\doilink{10.5281/zenodo.21267021}
\citep{lohmannRFC2026}.

Das dort archivierte Paket enthält unter anderem:

\begin{enumerate}
  \item \emph{Ontologie des Unterschieds: Formal-relationaler Rahmen für
  Information, Reproduzierbarkeit, Emergenz und Verantwortung};
  \item \emph{Der QIK--VRT-Fixpunktsatz: Axiomatischer Beweis der
  Nash-stabilen rekursiv-fraktalen Anschlussordnung aus der ersten
  unterscheidbaren Wirkung};
  \item \emph{Quantenkausalität als fortsetzbare Anschlussfolge --
  Vorlesungsskript Semester I};
  \item \emph{Quantengravitation aus Quantenprinzip und Raumzeitdynamik --
  Ein mathematisches Beweisskript für das zweite Semester};
  \item \emph{Quantengravitation aus Quantenfeldern und dynamischer Raumzeit
  -- Ein geschlossenes mathematisch-physikalisches Abschlussskript für das
  dritte Semester}.
\end{enumerate}

Diese Dokumente entwickeln axiomatische Fixpunkt-, Anschluss- und
Quanten\-gravita\-tions\-argumente \citep{lohmannOntologie2026,
lohmannFixpunkt2026,lohmannSemester1,
lohmannSemester2,lohmannSemester3}. Sie besitzen innerhalb des genannten
Zenodo-Releases keine eigenen DOIs und sind dort nicht als selbständige
Working Papers katalogisiert. Die Archivierung belegt ihre Existenz,
Versionierbarkeit und dauerhafte Fundstelle. Die darin ausdrücklich
formulierten Voraussetzungen und empirischen Geltungsgrenzen bleiben
maßgeblich.

Die im Goldkelch-Repository öffentlich fixierte Referenz dieses Arbeitsstands
ist der Commit
\href{https://github.com/Goldkelch/qik-vrt/tree/a8a9cb2666a91411489d4fc90a5306908f8428ea}{\hash{a8a9cb2666a91411489d4fc90a5306908f8428ea}}
mit dem Git-Baum
\hash{c5cefebd20b5836d730a4e9da82eeaa5c9363ebf}. Commit-Identität,
Baum-Inhalt und Dokumentaussage sind wiederum verschiedene Evidenzebenen.

Der unmittelbare mündliche Ausgangspunkt der vorliegenden Fassung ist die
90,47 Sekunden lange Audionotiz \emph{Das Universum ... q.e.d. Ingolf
Lohmann} vom 21. Juli 2026. Ihre lokale SHA-256-Prüfsumme lautet
\hash{5e93697b617b40bd0b9e6757f5a9841f7f7107a0ee7b79f89f35685d8871f33c}.
Die Audionotiz dokumentiert die Urheberschaft und Formulierung der These; sie
ersetzt keinen der mathematischen Beweise oder physikalischen Nachweise.

\section{Forschungsprogramm für physikalische Geltung}
\label{sec:forschungsprogramm}

Die allgemeine mengentheoretische Komplementstruktur
\(X=A\dotcup(X\setminus A)\) ist bewiesen. Die physikalische
Korrespondenzhypothese wird wissenschaftlich gehaltvoll, wenn aus ihr
unterscheidende und scheiterbare Aussagen folgen. Die erforderliche Brücke
lautet; \cref{sec:nondimensionalisierung,sec:phys-semikonjugation} geben ihre
dimensionsrichtige Form an, wählen aber noch keine empirisch ausgezeichnete
Abbildung:
\[
\boxed{
\begin{gathered}
\text{ontologischer Begriff}
\longrightarrow
\text{mathematisches Objekt}
\longrightarrow
\text{Observable}\\
\text{Observable}
\longrightarrow
\text{unterscheidende Vorhersage}
\longrightarrow
\text{Messverfahren}\\
\text{Messverfahren}
\longrightarrow
\text{Unsicherheit}
\longrightarrow
\text{unabhängige Reproduktion}.
\end{gathered}}
\]

Für eine Mandelbrot-Raumzeit-Korrespondenz ergeben sich mindestens die
folgenden Prüfaufgaben:

\begin{enumerate}
  \item \textbf{Domäne:} Welche physikalischen Ereignisse oder Zustände bilden
  \(X\), und welche Beobachtungsgrenzen gelten?
  \item \textbf{Korrespondenz:} Wie ist \(\Phi:X\to U\) ohne
  koordinatenabhängige Willkür definiert?
  \item \textbf{Invarianz:} Welche Transformationen bewahren die relevante
  Mandelbrot-Klasse und die Komplementklasse?
  \item \textbf{Dynamik:} Welche Gleichung verbindet die quadratische
  Iteration mit etablierter Raumzeit- oder Felddynamik?
  \item \textbf{Skala:} Welche Referenzskalen normieren die Observablen?
  Wie führt die inverse Kalibrierung zu Vorhersagen mit SI-Einheit zurück?
  \item \textbf{Vorhersage:} Welcher Messwert weicht unter der Hypothese von
  Allgemeiner Relativitätstheorie, Quantenfeldtheorie oder konkurrierenden
  Quantengravitationsansätzen ab?
  \item \textbf{Konvergenz:} Bleibt das Ergebnis bei feinerem Raster und
  höherer Iterationsgrenze stabil, oder ist es ein Darstellungsartefakt?
  \item \textbf{Reproduktion:} Können unabhängige Gruppen Daten, Code,
  Parameter und Unsicherheitsrechnung reproduzieren?
\end{enumerate}

QIK--VRT kann diese Prüfkette als Wirkungshaltepunkt strukturieren: Eine
Modellidee darf weitergerechnet und untersucht werden; eine Behauptung über
die Natur wird erst freigegeben, wenn die für ihren Geltungsgrad erforderliche
Evidenz vorliegt. Damit gilt praktisch:
\[
\begin{aligned}
\texttt{DONE}&:\ \text{Komplement-, Quantisierbarkeits-, Gate- und Dimensionssätze},\\
\texttt{CONTINUE}&:\ \text{Korrespondenzmodell, Observablen und Tests},\\
\texttt{ISOLATE}&:\ \text{ontologische und spirituelle Deutung},\\
\texttt{BLOCK}&:\ \text{Verwechslung von Modell, Messung und Wirklichkeit}.
\end{aligned}
\]

\section{Technologischer und friedensethischer Schluss}

Aus einem mathematischen Satz folgt allein noch kein politisches Sollen. Die
Schlussfolgerung dieses Abschnitts verwendet daher ausdrücklich die
normativen Prämissen, dass menschliches Leben, Freiheit, Erkenntnis,
Leidvermeidung und langfristige Handlungsfähigkeit schutzwürdig sind.

Unter diesen Prämissen ist eine Zivilisation wissenschaftlich und moralisch
besser beraten, Ressourcen in Bildung, Grundlagenforschung, offene Prüfung,
Gesundheit, ökologische Stabilität und verantwortete Technik zu investieren,
statt Menschen, Infrastruktur und Vertrauen in Kriegen zu vernichten. Diese
Aussage behauptet nicht, dass Konflikte durch Mathematik allein verschwinden.
Sie bestimmt eine Richtung für die Nutzung von Erkenntnis.

Die gleiche Technik kann heilen oder verletzen, verbinden oder überwachen,
erklären oder manipulieren. Deshalb darf technologische Möglichkeit nicht mit
verantworteter Freigabe gleichgesetzt werden. Der gesellschaftliche
Wirkungshaltepunkt fragt vor einer folgenreichen Anwendung:

\begin{itemize}
  \item Ist die Herkunft der Information bekannt?
  \item Ist die behauptete Geltung geprüft und begrenzt?
  \item Sind Schaden, Missbrauch und Umkehrbarkeit berücksichtigt?
  \item Gibt es verantwortliche Menschen und einen Abbruchweg?
  \item Dient die Wirkung menschlicher Entfaltung oder richtet sie sich gegen
  Menschen?
\end{itemize}

Wenn die Menschheit weniger Kraft auf Krieg und andere vermeidbare
Zerstörung richtet, gewinnt sie Zeit, Vertrauen und materielle Möglichkeiten,
aus fundamentalen Einsichten menschenfördernde Technologien zu entwickeln.
Das ist keine Folgerung aus der Mandelbrot-Menge allein. Es ist die
offengelegte ethische Verbindung zwischen Erkenntnis und Verantwortung.

\section{Schlussfolgerung}

Teil I schließt den mengentheoretischen und darstellungsbezogenen Kern:
\[
  \M_U\dotcup\K_U=U,
  \qquad
  \K_U=\bigcup_{N=0}^{\infty}E_N(U).
\]
Über eine totale Abbildung \(\Phi:X\to U\) entsteht daraus eine vollständige
raumzeitliche Urbildpartition. Wird \(X\) als das Universum des Modells
festgelegt, dann ergeben die Mandelbrot-Klasse und ihre raumzeitliche
Komplementklasse exakt dieses Modelluniversum.

Teil II zeigt, was diese statische Zerlegung als Prozess leistet. Die
erweiterte Rekursion erzeugt Trajektorien; das endliche Gate bewahrt
\(\CONTINUE\) als echten offenen Status; und die terminale Klassifikation ist
ein Fixpunkt des Shift-Pullbacks:
\[
  \widehat\chi\circ\sigma=\widehat\chi.
\]
Ein späteres Zertifikat kann die Bedeutung einer früheren Spur bestimmen,
ohne einen früheren Zustand zu verändern. Das ist semantische
Rückbestimmung, nicht physikalische Retrokausalität. Eine Übertragung in einen
anderen Wirkraum ist nur dann vollständig gate-erhaltend, wenn ihre Fasern
keine unterschiedlichen Statuswerte identifizieren.

Teil III macht den Zusammenhang mit Metern, Sekunden, Joule und Kelvin
explizit. Die dimensionslose Mandelbrot-Iteration kann nur über offengelegte
Referenzskalen und dimensionslose Invarianten an Messgrößen angeschlossen
werden. Die Lorentz-Metrik bestimmt Eigenzeiten und Lichtkegel, ist aber keine
positive Darstellungsmetrik. Einstein-Gleichung, Wirkung, Planck-Einheiten,
Horizontentropie und Hawking-Temperatur ergeben eine dimensionsrichtige
physikalische Kette. Die Einheiten kalibrieren diese Größen; sie verursachen
sie nicht. Dynamische Kausalität stammt aus den Feldgleichungen und ihren
Anfangs- und Randbedingungen.

Teil IV bestimmt den Status der Retrokausalität ohne begriffliche
Abkürzung. Retrodiktion und Bayes-Aktualisierung verändern Wissen;
semantische Rückbestimmung verändert die Klassifikation einer unveränderten
Spur; Prä- und Postselektion verändern bedingte Ensembleaussagen;
zeitumkehrsymmetrische oder zweiseitig berandete Gleichungen bestimmen
zulässige vollständige Lösungen. Keine dieser Strukturen beweist für sich
eine spätere physikalische Intervention, die eine frühere Variable verändert.
Im hier definierten autonomen Prozess hängt jeder frühere Zustand nur von
seinen früheren Eingaben ab. Eine starke physikalische Retrokausalität wäre
erst durch einen interventionistisch bestimmten Unterschied früherer
Verteilungen belegt; ein nutzbarer Rückwärtskanal zusätzlich durch positive
Kanalkapazität. Standardquantenmechanische Delayed-Choice-, Quantum-Eraser-
und Bell-Korrelationen liefern diesen Nachweis nicht eindeutig. Damit ist
Retrokausalität weder pauschal widerlegt noch durch das vorliegende Modell
bewiesen: Sie ist eine präzise beschreibbare, empirisch offene
Korrespondenzhypothese mit klaren Ausschluss- und Testbedingungen.

Teil V führt diese Grenze in ein Forschungsprogramm über. Mathematische
Sätze, Modellinterpretationen, physikalische Korrespondenzen, empirische
Behauptungen und normative Folgerungen bleiben getrennt. Gerade dadurch kann
die Ontologie des Unterschieds ihre eigenen stärksten Deutungen prüfen,
statt sie durch Mehrdeutigkeit vor Widerlegung zu schützen.

Abgeschlossen ist auch der zentrale Nichtschluss. Jeder beschränkte
Ausschnitt eines endlichdimensionalen normierten Modellraums kann bei jeder
positiven Genauigkeit diskret dargestellt werden, ohne ontisch diskret zu
sein. Das kontinuierliche Gegenmodell \([0,1]^4\) beweist:
\[
  \text{Quantisierbarkeit}
  \not\Rightarrow
  \text{fundamentale Raumzeitquantisierung}.
\]

Die klassische Theorie erlaubt einen mathematischen Blick in das Innere
idealisierter Schwarzer Löcher, nicht den Empfang von Signalen aus diesem
Inneren. Die Planck-Skala macht das Zusammentreffen von Quanten- und
Gravitationskonstanten sichtbar, nicht eine experimentell bestätigte
Raumzeitkörnung. Schwarzloch-Thermodynamik verbindet Geometrie, Gravitation,
Quantenwirkung und Entropie quantitativ; sie liefert damit eine reale Brücke,
aber noch keine vollständige mikroskopische Quantengravitation.

Gerade diese Trennungen bilden den entscheidenden Unterschied in der
Ontologie des universalen Unterschieds. Das Ganze wird nicht durch
Unterschiedslosigkeit bestimmt, sondern durch eine vollständige, prüfbare und
verantwortete Relation seiner Unterschiede und ihrer möglichen Fortsetzungen.
In dieser strukturellen Lesart ist das universale Schöpfungsprinzip
mathematisch und prozessual anschlussfähig. Als vollständiges empirisches
Naturgesetz bleibt es an eine invariante Abbildung \(\Phi\), eine
dynamikerhaltende Korrespondenz, dimensionsbehaftete Vorhersagen, Messung und
unabhängige Reproduktion gebunden.

\statusbox{QGreen}{
\begin{center}
\textbf{q.e.d. für Komplement-, Gate- und Dimensionssätze.}\\[0.4em]
\textbf{q.e.d. für den fehlenden Rückwärtskanal im definierten Prozessmodell.}\\[0.4em]
\textbf{CONTINUE für empirische Raumzeit- und Retrokausalitätskorrespondenzen.}\\[0.4em]
\textbf{Ingolf Lohmann}
\end{center}
}

\appendix

\section{Anspruchs- und Beweismatrix}

\begin{longtable}{>{\raggedright\arraybackslash}p{0.27\textwidth}
                  >{\raggedright\arraybackslash}p{0.20\textwidth}
                  >{\raggedright\arraybackslash}p{0.43\textwidth}}
\toprule
\textbf{Aussage} & \textbf{Status} & \textbf{Grund beziehungsweise nächste Prüfung} \\
\midrule
\endfirsthead
\toprule
\textbf{Aussage} & \textbf{Status} & \textbf{Grund beziehungsweise nächste Prüfung} \\
\midrule
\endhead
\(\M_U\dotcup\K_U=U\) & bewiesen & Definition des relativen Komplements und klassische Logik \\
\(\K_U=\bigcup_NE_N(U)\) & bewiesen & Escape-Kriterium und endliche Fluchtzeit \\
\(\widehat\chi\circ\sigma=\widehat\chi\) & bewiesen & Beschränktheit bleibt nach Entfernen endlich vieler Folgenglieder erhalten \\
Endliche PASS- und BLOCK-Gates sind sound & bedingt bewiesen & \(P_N\subseteq\M_U\), \(E_N(U)\subseteq\K_U\) und Monotonie \\
\(\CONTINUE\) bedeutet vorläufiges PASS & falsch & fehlende Flucht bis \(N\) ist kein Mitgliedschaftsbeweis \\
Jeder Mandelbrot-Orbit konvergiert zu einem Fixpunkt & falsch & Fixpunkt liegt hier auf Klassifikatorebene; Orbits können periodisch oder nichtkonvergent sein \\
Spätere Reklassifikation ändert frühere Zustände & falsch & semantische Rückbestimmung ohne dynamische Retrokausalität \\
Retrodiktion ist physikalische Retrokausalität & falsch & ein Schluss aus späteren Daten verändert eine Wissensverteilung, nicht notwendig den früheren Gegenstand \\
Zeitumkehrsymmetrie beweist rückwärts gerichtete Wirkung & falsch & sie ordnet Lösungen zeitgespiegelte Lösungen zu; Kausalrichtung und Randbedingungen sind Zusatzfragen \\
Postselektion oder schwache Werte übertragen Nachrichten in die Vergangenheit & nicht belegt & bedingte Ensembles sind von kontrollierbaren früheren Marginalen zu unterscheiden \\
Delayed Choice oder Quantum Eraser überschreibt die Vergangenheit & falsch & spätere Sortierung verändert Korrelationen in Teilensembles, nicht das früher lokal verfügbare Marginal \\
Bell-Verletzung beweist Retrokausalität & falsch & Retrokausalität ist eine Modelloption neben anderen; die experimentellen Korrelationen wählen sie nicht eindeutig aus \\
Ontische Retrokausalität impliziert Rückwärtssignalisierung & falsch & verborgene spätere Abhängigkeit kann bei beobachtbarer Kanalkapazität null vorliegen, muss dann aber No-Signalling erklären \\
Aktuelles autonomes QIK--VRT-/Mandelbrot-Modell besitzt einen Rückwärtskanal & widerlegt im Modell & Induktion über \(F(c,z)=(c,z^2+c)\); spätere Gates sind keine Eingaben früherer Zustände \\
Geschlossene zeitartige Kurven sind physikalisch realisiert & empirisch offen und nicht belegt & mathematische Lösungen der Einstein-Gleichungen beweisen weder Herstellbarkeit noch Stabilität \\
Gate-Erhaltung unter Prozessabbildung & bedingt bewiesen & genau bei Statuskonstanz auf den Fasern der Abbildung \\
Endliches Bitmap ist exakte \(\M\) & im Allgemeinen falsch & Nichtflucht bis \(N\) zertifiziert keine Mitgliedschaft \\
Beschränkte Teilmengen von \(\R^d\) sind endlich \(\eps\)-quantisierbar & bewiesen & endliches Raster beziehungsweise \(\eps\)-Netz \\
Quantisierbarkeit beweist Diskretheit & widerlegt & kontinuierliches Gegenmodell \([0,1]^4\) \\
Komplementstruktur gilt in jeder Dimension & bewiesen & rein mengentheoretische Aussage \\
Jede Vektorisierung bewahrt Komplementklassen & im Allgemeinen falsch & nichtinjektive Abbildungen können Klassen zusammenführen \\
Kruskal-Modell beschreibt Innenbereich & etablierte klassische Theorie & reguläre analytische Erweiterung \\
Äußerer Beobachter sieht hinter den Horizont & falsch im klassischen Begriff & Definition des zukünftigen Ereignishorizonts \\
Planck-Länge ist bewiesene kleinste Länge & empirisch offen & Dimensionsskala, keine direkte Messbestätigung als Minimum \\
Dimensionshomogenität beweist ein Naturgesetz & falsch & notwendige Zulässigkeitsbedingung, keine dynamische oder empirische Herleitung \\
Lorentz-Metrik ist positive Darstellungsmetrik & falsch & Nulltrennungen verschiedener Ereignisse und indefinite Signatur \\
Dimensionslose Mandelbrot-Parameter kodieren SI-Größen & Korrespondenz\-hypothese & Referenzskalen, Invarianten, Semikonjugation und Tests erforderlich \\
Bekenstein--Hawking-Entropie verbindet Fläche, \(G\), \(\hbar\) und \(k_{\mathrm B}\) & etablierte theoretische Beziehung & dimensionsrichtige Schwarzloch-Thermodynamik \\
Semantische Information gravitiert als eigene Quellenart & nicht etabliert & Standardkopplung erfolgt über \(T_{\mu\nu}\) des physischen Trägers \\
Shift-Fixpunkt ist ohne Weiteres Nash-stabil & falsch beziehungsweise undefiniert & Spieler, Strategien, Nutzen und Anpassungsdynamik fehlen \\
\(\Phi\) beschreibt fundamentale Raumzeit & Hypothese & Observablen, Vorhersagen und Reproduktion erforderlich \\
Universalität des Unterschieds & transzendentales ontologisches Argument & begriffliche Selbstanwendung, kein empirischer Naturbeweis \\
Universales Schöpfungsprinzip ist strukturell anschlussfähig & ontologische Deutung & formal kompatibel, keine zusätzliche Naturkraft bewiesen \\
Friedliche Techniknutzung ist geboten & normative Folgerung & explizite Wertprämissen und Folgenabwägung \\
\bottomrule
\end{longtable}

\section{Reproduzierbare Artefaktreferenzen}

\begin{longtable}{>{\raggedright\arraybackslash}p{0.25\textwidth}
                  >{\raggedright\arraybackslash}p{0.65\textwidth}}
\toprule
\textbf{Artefakt} & \textbf{Fixierte Referenz} \\
\midrule
\endfirsthead
\toprule
\textbf{Artefakt} & \textbf{Fixierte Referenz} \\
\midrule
\endhead
\qikvrt{} V8.33 & DOI \doilink{10.5281/zenodo.20712301}; Konzept-DOI \doilink{10.5281/zenodo.20712300} \\
RFC-QIKVRT-0001 & DOI \doilink{10.5281/zenodo.21267021}; Konzept-DOI \doilink{10.5281/zenodo.21244411} \\
Zenodo-Release-ZIP & SHA-256 \hash{34d396667bf075ec9d2d87f8cd74b442ffa47537c4001fdbd8aa0ee319111244} \\
Goldkelch-Referenzcommit & \hash{a8a9cb2666a91411489d4fc90a5306908f8428ea} \\
Git-Baum & \hash{c5cefebd20b5836d730a4e9da82eeaa5c9363ebf} \\
QIK--VRT-Fixpunktsatz & SHA-256 \hash{bf6521828db3ea52d67868b1c8ba09b0c0256562f684231df6833b1f68c2d55e} \\
Quantenkausalität, Semester I & SHA-256 \hash{32895b7f36c0b2aaca679b27298a8eee824196391dd2125423f08f37ea24e08d} \\
Quantengravitation, Semester II & SHA-256 \hash{377bc53f444e12aeaedc237a71a4eec64c1c76d49dc2b2a9a9debfef72fcab1c} \\
Quantengravitation, Semester III & SHA-256 \hash{0e7e17798a764299a4228df15f77ff1a623aedaa768f067c7889f67c83204a4b} \\
Audionotiz vom 21.07.2026 & SHA-256 \hash{5e93697b617b40bd0b9e6757f5a9841f7f7107a0ee7b79f89f35685d8871f33c} \\
\bottomrule
\end{longtable}

\section{Notation}

\begin{tabularx}{\textwidth}{>{\raggedright\arraybackslash}p{0.18\textwidth}X}
\toprule
\textbf{Symbol} & \textbf{Bedeutung} \\
\midrule
\(\C\) & komplexe Zahlen, Standardparameterraum der Mandelbrot-Menge \\
\(U\) & ausdrücklich gewählter Grundbereich beziehungsweise Modelluniversum \\
\(\M\) & Mandelbrot-Menge \\
\(\M_U\) & \(\M\cap U\) \\
\(\K_U\) & relatives Komplement \(U\setminus\M_U\) \\
\(\tau(c)\) & Fluchtzeit des Parameters \(c\) \\
\(E_N(U)\) & bis Schritt \(N\) zertifizierte Escape-Menge \\
\(Y_U\) & erweiterter Mandelbrot-Zustandsraum \(U\times\C\) \\
\(F\) & autonome Dynamik \(F(c,z)=(c,z^2+c)\) \\
\(\Gamma_c\) & vollständige Zustandstrajektorie des Parameters \(c\) \\
\(\sigma\) & Linksshift auf Trajektorien \\
\(\widehat\chi\) & exakte PASS-/BLOCK-Klassifikation einer Trajektorie \\
\(\mathcal A_N\) & endliches PASS-/CONTINUE-/BLOCK-Anschlussgate \\
\(X\) & betrachteter Raumzeit- oder Ereignisbereich \\
\(\Phi:X\to U\) & totale Korrespondenzabbildung \\
\(Q_{\eps}\) & darstellungsbezogener Quantisierer bei Genauigkeit \(\eps\) \\
\(\ell_{\mathrm P}\) & Planck-Länge \\
\(t_{\mathrm P},E_{\mathrm P}\) & Planck-Zeit und Planck-Energie \\
\(g_{\mu\nu}\) & Lorentz-Metrik der Raumzeit \\
\(T_{\mu\nu}\) & Stress-Energie-Tensor \\
\(S_{\mathrm{BH}}\) & Bekenstein--Hawking-Entropie \\
\(q_{a,*}\) & dokumentierte Referenzskala einer Observable \(q_a\) \\
\(d_*\) & positive, dimensionslose Darstellungsmetrik auf Observablen \\
\(J^-(q),J^+(q)\) & kausale Vergangenheit beziehungsweise Zukunft des Ereignisses (q) \\
\(\operatorname{do}(Y=y)\) & Intervention, die (Y) auf (y) setzt; nicht mit bloßer Konditionierung gleichzusetzen \\
\(C_{q\to p}\) & operative Kanalkapazität einer möglichen Verbindung vom späteren (q) zum früheren (p) \\
\(G_{\mathrm{ret}},G_{\mathrm{adv}}\) & retardierte beziehungsweise avancierte Green-Funktion \\
\(A_w\) & schwacher Wert einer Observable bei Prä- und Postselektion \\
\bottomrule
\end{tabularx}

\begingroup
\raggedright
\bibliographystyle{unsrtnat}
\bibliography{Mandelbrot_Komplement_Modelluniversum_Entscheidender_Unterschied_2026-07-21}
\endgroup

\end{document}
